\documentclass[12pt,openany]{report}

% ==== Paquetes necesarios ====
\usepackage{tikz}
\usepackage{afterpage}
\usepackage[most]{tcolorbox}
\usepackage{graphicx}
\newcounter{pagebeforecert}
\usepackage{xcolor}
\usepackage{fontawesome5}
\usetikzlibrary{arrows.meta,positioning,shapes.geometric,fit, calc}

\usepackage{array}
\usepackage{tabularx}
\usepackage{ragged2e}
\usepackage{xcolor}
\usepackage{colortbl}
\usepackage{adjustbox}
\usepackage{ragged2e}
\usepackage{parskip}

\usepackage[utf8]{inputenc}
\usepackage[T1]{fontenc}
\usepackage[spanish, es-nodecimaldot]{babel}
\usepackage{helvet}
\renewcommand{\familydefault}{\sfdefault}
\usepackage{microtype}
\usepackage{geometry}
\geometry{letterpaper, top=4.2cm, bottom=2.1cm, left=2.1cm, right=2.1cm, headheight=60pt, headsep=1.5cm}
\usepackage{multirow}
\usepackage{fancyhdr}
\usepackage{setspace}
\onehalfspacing


\usepackage[hidelinks]{hyperref}
\hypersetup{
  pdftitle={Manual de Seguridad Operacional Skyfilms},
  pdfauthor={Skyfilms},
  pdfsubject={RAC 100 - Manual de Seguridad Operacional},
  pdfcreator={LaTeX},
  pdfproducer={Overleaf}
}
\usepackage{graphicx}

\usepackage{float}
\usepackage{longtable, booktabs}
\usepackage{array}
\usepackage{enumitem}
\usepackage{csquotes}
\usepackage[acronym]{glossaries}
\makeglossaries
% --- Acrónimos y abreviaturas ---

\newacronym{aaae}{AAAE}{Autoridad Aeronáutica de la Aviación de Estado}
\newacronym{aac}{AAC}{Autoridad de Aeronáutica Civil}
\newacronym{agl}{AGL}{Altura sobre el nivel del terreno}
\newacronym{arp}{ARP}{Punto de referencia del aeródromo}
\newacronym{atc}{ATC}{Servicio de control de tránsito aéreo}
\newacronym{ats}{ATS}{Servicios de tránsito aéreo}
\newacronym{bvlos}{BVLOS}{Beyond Visual Line of Sight}
\newacronym{cdm}{CDM}{Toma de decisiones en colaboración}
\newacronym{ceac}{CEAC}{Centro de entrenamiento de aeronáutica civil}
\newacronym{ciac}{CIAC}{Centro de instrucción de aeronáutica civil}
\newacronym{conops}{CONOPS}{Concepto de Operaciones}
\newacronym{dona}{DONA}{Dirección de Operaciones de Navegación Aérea}
\newacronym{evlos}{EVLOS}{Visibilidad en línea de vista extendida}
\newacronym{fato}{FATO}{Área de aproximación final y despegue}
\newacronym{fhss}{FHSS}{Frequency Hopping Spread Spectrum}
\newacronym{fpv}{FPV}{First Person View}
\newacronym{gnss}{GNSS}{Sistema Global de Navegación por Satélite}
\newacronym{metar}{METAR}{Meteorological Aerodrome Report}
\newacronym{mo}{MO}{Manual de Operaciones}
\newacronym{notam}{NOTAM}{Notice to Airmen}
\newacronym{opsp}{OpSpecs}{Especificaciones Operacionales}
\newacronym{osd}{OSD}{Telemetría en pantalla}
\newacronym{pmbo}{PMBO}{Peso máximo bruto de operación}
\newacronym{pmu}{PMU}{Puesto de mando unificado}
\newacronym{racae}{RACAE}{Reglamento Aeronáutico Colombiano de la Aviación de Estado}
\newacronym{rca}{RAC}{Reglamentos Aeronáuticos de Colombia}
\newacronym{rce}{RCE}{Responsabilidad Civil Extracontractual}
\newacronym{rpa}{RPA}{Aeronave pilotada a distancia}
\newacronym{rpas}{RPAS}{Sistema de aeronave pilotada a distancia}
\newacronym{sa}{SA}{Área de seguridad}
\newacronym{sarps}{SARPS}{Normas y métodos recomendados publicados por la OACI}
\newacronym{sms}{SMS}{Sistema de Gestión de la Seguridad Operacional}
\newacronym{taf}{TAF}{Terminal Aerodrome Forecast (Pronóstico de Aeródromo)}
\newacronym{tlof}{TLOF}{Área de contacto y elevación.}
\newacronym{ua}{UA}{Aeronave no tripulada}
\newacronym{uaeac}{UAEAC}{Unidad Administrativa Especial de Aeronáutica Civil}
\newacronym{uas}{UAS}{Sistema de Aeronaves no Tripuladas}
\newacronym{vlos}{VLOS}{Visual Line of Sight}
\newacronym{vmc}{VMC}{Condiciones meteorológicas visuales}
\newacronym{vtol}{VTOL}{Vertical Take-Off and Landing / Despegue y aterrizaje vertical}
\newacronym{znvd}{ZNVD}{Zona de No Vuelo Dron}

\usepackage[backend=biber,style=ieee]{biblatex}
\addbibresource{referencias.bib}

% ======================================================================
% AJUSTE NATIVO DEL ESPACIO EN INICIOS DE CAPÍTULO (SIN ERRORES)
% ======================================================================
\makeatletter
% Modificamos la macro original de LaTeX que inserta el espacio en blanco gigante
\def\@makechapterhead#1{%
  % El valor original es \vspace*{50\p@}. Lo cambiamos a 0 o negativo para subirlo
  \vspace*{-20\p@}% <--- ESTE ES EL NÚMERO QUE PUEDES CAMBIAR (ej: -10, -30)
  {\parindent \z@ \raggedright \normalfont
    \interlinepenalty\@M
    \Huge \bfseries #1\par\nobreak
    \vskip 40\p@
  }}
  
% Y hacemos lo mismo para los capítulos sin número (\chapter*) como los Apéndices
\def\@makeschapterhead#1{%
  \vspace*{-20\p@}% <--- EL MISMO VALOR QUE EL ANTERIOR PARA QUE QUEDEN IGUALES
  {\parindent \z@ \raggedright
    \normalfont
    \interlinepenalty\@M
    \Huge \bfseries  #1\par\nobreak
    \vskip 40\p@
  }}
\makeatletter
% ======================================================================


% ======================================================================
% CONFIGURACIÓN DEL ENCABEZADO SMS (A PRUEBA DE FALLOS)
% ======================================================================
\pagestyle{fancy}
\fancyhf{}
\renewcommand{\headrulewidth}{0pt}

\fancyhead[C]{
    \renewcommand{\arraystretch}{1.6} % Da aire vertical a las celdas
    \begin{tabular}{|m{3.8cm}|m{6.5cm}|m{4.8cm}|}
        \hline
        % Al poner el ancho fijo (3.8cm y 6.5cm) en el multirow, LaTeX centra todo automáticamente
        \multirow{3}{3.8cm}{\centering \includegraphics[width=3.2cm]{Imagenes/LOGO.png}} & 
        \multirow{3}{6.5cm}{\centering \\[0.2]\textbf{MANUAL DE SEGURIDAD}\\[0.1cm] \textbf{OPERACIONAL SMS}} & 
        \textbf{Versión: 1.0} \\
        \cline{3-3}
        & & \textbf{Código: SMS-SKY-001} \\
        \cline{3-3}
        & & \textbf{Fecha de Emisión: 23/06/2026} \\
        \hline
    \end{tabular}
}
\fancyfoot[R]{\thepage}

% Copia exacta para las páginas donde inicia un capítulo
\fancypagestyle{plain}{
    \fancyhf{}
    \renewcommand{\headrulewidth}{0pt}
    \fancyhead[C]{
        \renewcommand{\arraystretch}{1.6}
        \begin{tabular}{|m{3.8cm}|m{6.5cm}|m{4.8cm}|}
            \hline
            \multirow{3}{3.8cm}{\centering \includegraphics[width=3.2cm]{Imagenes/LOGO.png}} & 
            \multirow{3}{6.5cm}{\centering \\[0.1]\textbf{MANUAL DE SEGURIDAD}\\[0.1cm] \textbf{OPERACIONAL SMS}} & 
            Versión: 01 \\
            \cline{3-3}
            & & Código: SMS-SKY-001 \\
            \cline{3-3}
            & & Fecha de Emisión: 23/06/2026 \\
            \hline
        \end{tabular}
    }
    \fancyfoot[R]{\thepage}
}
% =========================================================================

\begin{document}
\begin{titlepage}
    \centering
    {\Large \textbf{SKYFILMS}}\\[6pt]
    {\huge \textbf{Manual de Seguridad Operacional}}\\[6pt]
    {\Large \textbf{SMS}}\\[6pt]
    {\large Cumplimiento RAC 100}\\[30pt]
    \includegraphics[width=0.6\textwidth]{Imagenes/LOGO.png}\\[16pt]
    \begin{tabular}{ll}
    \\
    \\
    \\
    \\
    \\
    \\
    
        \textbf{Código}: & SMS-SKY-001\\
        \textbf{Versión}: & 01\\
            
    \end{tabular}
    \vfill
\end{titlepage}



\chapter*{Lista de Páginas Efectivas}
\addcontentsline{toc}{chapter}{Lista de Páginas Efectivas}
\begin{center}
\begin{longtable}{|c|p{7.5cm}|c|c|c|}
\hline
\rowcolor[HTML]{D9D9D9} 
\textbf{Cap./Sección} & \textbf{Título} & \textbf{Páginas} & \textbf{Revisión} & \textbf{Fecha} \\ \hline
\endfirsthead

\hline
\rowcolor[HTML]{D9D9D9} 
\textbf{Cap./Sección} & \textbf{Título} & \textbf{Páginas} & \textbf{Revisión} & \textbf{Fecha} \\ \hline
\endhead

\hline
\endfoot

\hline
\endlastfoot

Preliminares & Lista de Páginas Efectivas & 1 & 01 & 02/06/2026 \\ \hline
Preliminares & Control Manual Enmiendas y Revisiones & 2 - 4 & 01 & 02/06/2026 \\ \hline
Preliminares & Lista de Distribución & 5 & 01 & 02/06/2026 \\ \hline
Preliminares & Definiciones y Acrónimos & 6 - 9 & 01 & 02/06/2026 \\ \hline
Preliminares & Índice General & 10 - 12 & 01 & 02/06/2026 \\ \hline
1 & Introducción & 13 - 15 & 01 & 02/06/2026 \\ \hline
2 & Política y Objetivos & 16 - 37 & 01 & 02/06/2026 \\ \hline
3 & Gestión de Riesgos de Seguridad Operacional & 38 - 47 & 01 & 02/06/2026 \\ \hline
4 & Aseguramiento de la Seguridad Operacional & 48 - 54 & 01 & 02/06/2026 \\ \hline
5 & Promoción de la Seguridad Operacional & 55 - 59 & 01 & 02/06/2026 \\ \hline
Apéndices & Anexo 1: Objetivos de Seguridad KPI's & 61+ & 01 & 02/06/2026 \\ \hline
Apéndices & Anexo 2: Contactos de Emergencia & 61+ & 01 & 02/06/2026 \\ \hline
Apéndices & Anexo 3: Proc. Control Documentos y Registros & 61+ & 01 & 02/06/2026 \\ \hline
Apéndices & Anexo 4: Listado Reportes Obligatorios M.O.R. & 61+ & 01 & 02/06/2026 \\ \hline
Apéndices & Anexo 5: Guía Formato Notificación Peligros & 61+ & 01 & 02/06/2026 \\ \hline
Apéndices & Anexo 6: Guía Evaluación de Riesgos FRAT & 61+ & 01 & 02/06/2026 \\ \hline
Apéndices & Anexo 7: Registro de Peligros & 61+ & 01 & 02/06/2026 \\ \hline
Apéndices & Anexo 8: Matriz de Evaluación (Reactivo) & 61+ & 01 & 02/06/2026 \\ \hline
Apéndices & Anexo 9: Indicadores de Rendimiento SPI's & 61+ & 01 & 02/06/2026 \\ \hline
Apéndices & Anexo 10: Proc. de Gestión del Cambio & 61+ & 01 & 02/06/2026 \\ \hline
Apéndices & Anexo 11: Plan de Mejora Continua & 61+ & 01 & 02/06/2026 \\ \hline
Apéndices & Anexo 12: Programa de Instrucción SMS & 61+ & 01 & 02/06/2026 \\ \hline
Apéndices & Documentación de Referencia & 61+ & 01 & 02/06/2026 \\ \hline

\end{longtable}
\end{center}


\chapter*{Control Manual Enmiendas y Revisiones}
\addcontentsline{toc}{chapter}{Control Manual Enmiendas y Revisiones}
%%%%%%%%%%%%%%%%%%%%%%%%%%%%%%%%%%%%%%%%%%%%%%%%%%%%%%%%%%%%%%%%%%%%%%%%%%%%%%%%
% CONTROL DE ENMIENDAS Y REVISIONES
%%%%%%%%%%%%%%%%%%%%%%%%%%%%%%%%%%%%%%%%%%%%%%%%%%%%%%%%%%%%%%%%%%%%%%%%%%%%%%%%

\addcontentsline{toc}{chapter}{Control de Enmiendas y Revisiones}

El presente Manual de Seguridad Operacional (SMS) es un documento controlado del sistema documental de Skyfilms S.A.S. Toda modificación realizada al documento deberá ser identificada, aprobada, registrada y distribuida conforme al procedimiento de control documental establecido por la organización.

Las modificaciones al Manual se clasifican de la siguiente manera:

\begin{itemize}
    \item \textbf{Revisión:} modificación menor que no altera el alcance, la estructura o las condiciones de aprobación del Manual de Seguridad Operacional (SMS).
    
    \item \textbf{Enmienda:} modificación mayor que afecta procedimientos operacionales, responsabilidades, estructura documental, flota UAS, política, objetivos o indicadores de seguridad operacional.
\end{itemize}

Toda revisión o enmienda deberá ser aprobada por el Gerente General y registrada en las tablas siguientes con el fin de mantener la trazabilidad histórica del documento.

\vspace{0.7cm}

%%%%%%%%%%%%%%%%%%%%%%%%%%%%%%%%%%%%%%%%%%%%%%%%%%%%%%%%%%%%%%%%%%%%%%%%%%%%%%%%
% REGISTRO DE REVISIONES
%%%%%%%%%%%%%%%%%%%%%%%%%%%%%%%%%%%%%%%%%%%%%%%%%%%%%%%%%%%%%%%%%%%%%%%%%%%%%%%%

\section*{Registro de Revisiones}

\renewcommand{\arraystretch}{1.35}

\begin{longtable}{|>{\centering\arraybackslash}p{1.2cm}|>{\centering\arraybackslash}p{2.1cm}|p{4.2cm}|>{\centering\arraybackslash}p{2.2cm}|>{\centering\arraybackslash}p{1.8cm}|p{2.5cm}|}

\hline

\rowcolor{gray!20}

\textbf{Rev.} &
\textbf{Sección} &
\textbf{Descripción del cambio} &
\textbf{Páginas afectadas} &
\textbf{Fecha} &
\textbf{Responsable}

\\
\hline
\endfirsthead

\hline

\rowcolor{gray!20}

\textbf{Rev.} &
\textbf{Sección} &
\textbf{Descripción del cambio} &
\textbf{Páginas afectadas} &
\textbf{Fecha} &
\textbf{Responsable}

\\
\hline

\endhead

0 &
Emisión inicial &
Emisión original del Manual de Seguridad Operacional (SMS) &
Todas &
23/06/2026 &
Gerente General

\\
\hline

&
&
&
&
&
\\
\hline

&
&
&
&
&
\\
\hline

&
&
&
&
&
\\
\hline

&
&
&
&
&
\\
\hline

\end{longtable}

\vspace{0.7cm}

%%%%%%%%%%%%%%%%%%%%%%%%%%%%%%%%%%%%%%%%%%%%%%%%%%%%%%%%%%%%%%%%%%%%%%%%%%%%%%%%
% HISTORIAL DE VERSIONES
%%%%%%%%%%%%%%%%%%%%%%%%%%%%%%%%%%%%%%%%%%%%%%%%%%%%%%%%%%%%%%%%%%%%%%%%%%%%%%%%

\section*{Historial de Versiones}

\renewcommand{\arraystretch}{1.35}

\begin{longtable}{|>{\centering\arraybackslash}p{1.6cm}|p{3.0cm}|p{5.6cm}|>{\centering\arraybackslash}p{1.9cm}|p{2.3cm}|}

\hline

\rowcolor{gray!20}

\textbf{Versión} &
\textbf{Secciones afectadas} &
\textbf{Descripción de la enmienda} &
\textbf{Fecha} &
\textbf{Aprobó}

\\
\hline

\endfirsthead

\hline

\rowcolor{gray!20}

\textbf{Versión} &
\textbf{Secciones afectadas} &
\textbf{Descripción de la enmienda} &
\textbf{Fecha} &
\textbf{Aprobó}

\\
\hline

\endhead

1.0 &
Documento completo &
Primera emisión del Manual de Seguridad Operacional (SMS). &
23/06/2026 &
Gerente General

\\
\hline

&
&
&
&
\\
\hline

&
&
&
&
\\
\hline

&
&
&
&
\\
\hline

&
&
&
&
\\
\hline

\end{longtable}

\vspace{0.4cm}

\begin{tcolorbox}[
colback=gray!5,
colframe=black,
title=Control Documental]

La versión vigente del Manual de Seguridad Operacional (SMS) corresponde a la última versión aprobada por el Gerente General de Skyfilms S.A.S. Toda copia anterior deberá retirarse de circulación una vez emitida una nueva versión del documento.

\end{tcolorbox}



\chapter*{Lista de Distribución}
\addcontentsline{toc}{chapter}{Lista de Distribución}
\begin{longtable}{@{}lp{5cm}l@{}}
\toprule
\textbf{Rol} & \textbf{Formato} \\
\midrule
Gerencia General & PDF (Editable) \\
Gerente de Seguridad Operacional  & PDF (Editable) \\
Jefe de Pilotos  & PDF  \\
Pilotos UAS & PDF  \\
Técnico UAS & PDF\\
Aseguradora / Cliente (cuando aplique)  & PDF  \\
\bottomrule
\end{longtable}


\chapter*{Definiciones y Acrónimos}
\addcontentsline{toc}{chapter}{Definiciones y Acrónimos}
El presente capítulo establece la terminología oficial y el significado de las siglas utilizadas a lo largo de este Manual de Gestión de Seguridad Operacional (SMS). Su objetivo es garantizar una interpretación uniforme, precisa y estandarizada de los conceptos por parte de todo el personal de SKYFILMS, la Aerocivil y terceros involucrados en la operación.\\

\textbf{1. Acrónimos y Siglas}

Para los fines del presente manual, se aplicarán las siguientes siglas y acrónimos:

\begin{longtable}{|>{\centering\arraybackslash}m{0.17\linewidth}|m{0.6\linewidth}|}
\hline
\rowcolor{gray!20}
\textbf{Sigla / Acrónimo} & \textbf{Significado (Inglés / Español)} \\
\hline
\endfirsthead % Indica que este encabezado solo se muestra en la primera página de la tabla

\hline
\rowcolor{gray!20}
\textbf{Sigla / Acrónimo} & \textbf{Significado (Inglés / Español)} \\
\hline
\endhead % Indica que este encabezado se repetirá automáticamente si la tabla pasa a otra hoja

AOC / CDO & Certificado de Explotador de Servicios Aéreos / Certificado de Operador. \\ \hline
C2 & Command and Control / Enlace de Mando y Control del UAS. \\ \hline
ERP / PRE & Emergency Response Plan / Plan de Respuesta ante Emergencias. \\ \hline
FRAT & Flight Risk Assessment Tool / Herramienta de Evaluación de Riesgos de Vuelo. \\ \hline
KPI & Key Performance Indicator / Indicador Clave de Rendimiento. \\ \hline
MAC & Mid-Air Collision / Colisión en el Aire. \\ \hline
MIE & Matriz de Identificación y Evaluación de Riesgos. \\ \hline
MOR & Mandatory Occurrence Report / Reporte Obligatorio de Eventos de Seguridad Operacional. \\ \hline
OACI / ICAO & Organización de Aviación Civil Internacional / International Civil Aviation Organization. \\ \hline
RAC & Reglamentos Aeronáuticos de Colombia. \\ \hline
RPA & Remotely Piloted Aircraft / Aeronave Pilotada a Distancia. \\ \hline
SA & Safety Assurance / Aseguramiento de la Seguridad Operacional. \\ \hline
SMS & Safety Management System / Sistema de Gestión de Seguridad Operacional. \\ \hline
SPI & Safety Performance Indicator / Indicador de Rendimiento de Seguridad Operacional. \\ \hline
SPT & Safety Performance Target / Meta de Rendimiento de Seguridad Operacional. \\ \hline
SRM & Safety Risk Management / Gestión de Riesgos de Seguridad Operacional. \\ \hline
TRE & Tecnología, Reglamento y Entrenamiento (Metodología de mitigación de riesgos RAC 100). \\ \hline
Aerocivil & Unidad Administrativa Especial de Aeronáutica Civil. \\ \hline
UAS & Unmanned Aircraft System / Sistema de Aeronave No Tripulada. \\ \hline
VLOS & Visual Line of Sight / Operación en Línea de Vista Visual. \\ \hline
DIACC & Grupo de Investigación de Accidentes de la Aerocivil. \\ \hline
FFHH & Factores Humanos. \\ \hline
SAG & Safety Action Group / Grupo de Acción de Seguridad Operacional. \\ \hline
SRB & Safety Review Board / Junta de Revisión de Seguridad Operacional. \\ \hline
SDCPS & Safety Data Collection and Processing Systems / Sistemas de Recopilación y Procesamiento de Datos de Seguridad Operacional. \\ \hline
VOR & Voluntary Occurrence Report / Reporte Voluntario. \\ \hline
DIS & Documento de Información Suplementaria entregado a la Aerocivil. \\ \hline
OpSpecs & Especificaciones de Operación. \\ \hline
MO & Manual de Operaciones de Skyfilms. \\ \hline
MCM & Manual de Control de Mantenimiento de Skyfilms. \\ \hline
RUAS & Registro de Unidades de Aeronaves no Tripuladas. \\ \hline
FPV & First Person View / Vista en Primera Persona. \\ \hline
EPP & Elementos de Protección Personal. \\ \hline
SST & Seguridad y Salud en el Trabajo. \\ \hline

% Nota: En longtable la leyenda (caption) y etiqueta (label) se ponen adentro, al final de la tabla
\caption{Acrónimos del Manual SMS}
\label{tab:Acronimos}
\end{longtable}

\textbf{2. Definiciones}

Con el fin de asegurar una comprensión integral de los procesos de SKYFILMS, se adoptan las siguientes definiciones oficiales (alineadas con el RAC 1 y RAC 219):

\begin{enumerate}
    \item \textbf{Aceptabilidad del Riesgo:} El nivel de riesgo que una organización está dispuesta a asumir tras haber implementado las defensas y controles correspondientes, basándose en la matriz de tolerabilidad.

    \item \textbf{Auditoría de Seguridad Operacional:} Examen sistemático, independiente y documentado para determinar si las actividades y los resultados del SMS cumplen con las disposiciones planificadas y la normatividad aeronáutica.

    \item \textbf{Consecuencia:} El resultado potencial, impacto o efecto de que un peligro se materialice (Ej. Daños a terceros, lesiones, pérdida del equipo).

    \item \textbf{Cultura Justa (Just Culture):} Un entorno donde no se castiga a los empleados por acciones, omisiones o decisiones tomadas de acuerdo con su experiencia y capacitación, pero en el cual no se toleran la negligencia grave, las infracciones intencionales ni los actos destructivos.

    \item \textbf{Ejecutivo Responsable:} Persona única en la organización (SKYFILMS) que actúa al más alto nivel directivo y cuyas responsabilidades, basadas en su autoridad y atribuciones financieras, son absoluta y totalmente indelegables frente a la eficacia del SMS.

    \item \textbf{Gestión del Cambio:} Proceso formal para identificar, evaluar y mitigar los riesgos de seguridad operacional que surgen como resultado de modificaciones en la estructura de la empresa, adquisición de nuevos equipos (UAS), o cambios normativos.

    \item \textbf{Mejora Continua:} Proceso recurrente que permite incrementar la eficacia y eficiencia del Sistema de Gestión de Seguridad Operacional mediante la evaluación de indicadores (SPIs), auditorías y lecciones aprendidas.

    \item \textbf{Mitigación (Defensas):} Medida, control o acción correctiva adoptada para reducir la probabilidad o la severidad de los riesgos asociados a un peligro identificado. En SKYFILMS, se categorizan bajo el modelo TRE (Tecnología, Reglamento, Entrenamiento).

    \item \textbf{Peligro:} Condición, objeto o actividad que podría provocar un incidente o accidente de aviación, o contribuir al mismo.

    \item \textbf{Probabilidad:} La posibilidad de que ocurra un evento inseguro o una consecuencia indeseada al interactuar con un peligro.

    \item \textbf{Rendimiento en Materia de Seguridad Operacional:} El logro de una organización con respecto a la seguridad operacional, definido mediante la medición de sus indicadores (SPI) y metas (SPT).

    \item \textbf{Riesgo de Seguridad Operacional:} La evaluación combinada de la probabilidad y la severidad proyectadas de las consecuencias de un peligro.

    \item \textbf{Seguridad Operacional (Safety):} El estado en el cual los riesgos asociados a las actividades de aviación (operación de UAS) se reducen y controlan a un nivel aceptable.

    \item \textbf{Severidad:} La magnitud o gravedad del impacto de las consecuencias de un peligro en caso de que este se materialice.

    \item \textbf{Sistema de Aeronave No Tripulada (UAS):} Una aeronave y sus elementos conexos, que se operan sin piloto a bordo (incluye la estación de control, el enlace C2 y los equipos de apoyo).

    \item \textbf{Gerente SMS:} Persona designada por el Gerente General para administrar el SMS. Es el único nombre de este cargo en todo el Manual (equivale a Asesor SMS o Gerente de Seguridad Operacional en otros documentos).

    \item \textbf{Junta de Revisión de Seguridad Operacional (SRB):} Grupo directivo, presidido por el Gerente General, que revisa el desempeño del SMS y decide sobre recursos y riesgos de nivel alto.

    \item \textbf{Grupo de Acción de Seguridad Operacional (SAG):} Grupo operativo, dirigido por el Gerente SMS, que analiza reportes, riesgos y acciones y las lleva a cabo.

    \item \textbf{Riesgo Residual:} Riesgo que queda después de aplicar los controles.

    \item \textbf{Riesgo Agregado:} Suma de varios riesgos pequeños que, juntos, pueden ser mayores que cada uno por separado.
\end{enumerate}



\tableofcontents
\clearpage


% Capítulo 1: Información General
\chapter{Introducción}
\label{chap:introduccion}

\textbf{SKYFILMS}, una empresa colombiana líder en producción audiovisual con drones, se compromete a brindar servicios de alta calidad de manera segura y responsable. Con el objetivo de garantizar la seguridad de nuestras operaciones, hemos desarrollado este Manual del Sistema de Gestión de Seguridad Operacional (SMS), alineado con los requisitos establecidos en el Reglamento Aeronáutico Civil (RAC) en sus versiones 100 y 219.\\

Este manual establece las directrices para la gestión sistemática de la seguridad operacional, documentando las estructuras organizacionales, líneas de responsabilidad, políticas y procedimientos necesarios para mantener nuestras operaciones bajo control. Para \textbf{SKYFILMS}, el SMS no es un componente aislado, sino un enfoque de negocio fundamental que se integra en todos los procesos de planificación, ejecución y evaluación de nuestras misiones audiovisuales.\\

A través de este documento, la organización consolida los cuatro pilares fundamentales del sistema: la Política y Objetivos de Seguridad Operacional, la Gestión de Riesgos, el Aseguramiento de la Seguridad y la Promoción de la Seguridad Operacional. Nuestro propósito es transitar de un cumplimiento normativo básico hacia una gestión proactiva y predictiva, garantizando la identificación temprana de peligros y la implementación de defensas eficaces que mantengan los riesgos en un nivel tan bajo como sea razonablemente práctico.


\section{Alcance}

El presente manual establece el alcance del Sistema de Gestión de Seguridad Operacional (SMS) de \textbf{SKYFILMS}, el cual es de estricto cumplimiento para todo el personal y aplicable a todas las actividades relacionadas con la prestación de nuestros servicios aéreos.\\

\textbf{Alcance Operacional:}

El SMS aplica a todas las operaciones aéreas realizadas por \textbf{SKYFILMS} utilizando Sistemas de Aeronaves No Tripuladas (UAS). Esto incluye las fases de planificación, preparación, ejecución y post-vuelo. Nuestro ámbito de operación abarca entornos diversos y complejos, incluyendo zonas urbanas, áreas rurales y operaciones sobre multitudes, empleando una variedad de aeronaves no tripuladas según las autorizaciones correspondientes.\\

\textbf{Alcance Organizacional:}\\
La seguridad operacional es una responsabilidad transversal en \textbf{SKYFILMS}. Por lo tanto, los lineamientos, políticas y procedimientos descritos en este manual aplican a todos los niveles estructurales de la organización, vinculando expresamente a:\\

\begin{itemize}
    \item \textbf{Gerencia General} (como Ejecutivo Responsable de la provisión de recursos y del SMS).

    \item \textbf{Gerente de Seguridad Operacional} (responsable de la administración del sistema).

    \item Jefe de Pilotos.

    \item Pilotos al mando, observadores y cualquier personal técnico o de apoyo involucrado en la operación.
\end{itemize}


\textbf{Aplicabilidad:}\\
Para garantizar un enfoque sistémico, los componentes del SMS se integran de manera obligatoria en los siguientes procesos de la organización:\\

\begin{itemize}
    \item \textbf{Operaciones:} Planificación, supervisión, autorización y ejecución segura de todos los vuelos UAS.

    \item \textbf{Mantenimiento:} Procesos de inspección, verificación pre-vuelo, conservación y gestión de la aeronavegabilidad continua de los equipos y estaciones de control.

    \item \textbf{Gestión de Riesgos:} Identificación proactiva y predictiva de peligros, evaluación de la probabilidad y severidad, y aplicación de medidas de mitigación.

    \item \textbf{Aseguramiento de la Seguridad:} Reporte voluntario y obligatorio de eventos de seguridad, investigación, análisis de tendencias y auditorías internas para la mejora continua.

    \item \textbf{Promoción de la Seguridad:} Capacitación continua del personal, evaluación de competencias y promoción activa de la comunicación y una cultura justa.
\end{itemize}


\section{Objetivo del Manual}
El objetivo principal de este manual es establecer un marco de trabajo que permita identificar, evaluar y mitigar los riesgos asociados a nuestras operaciones con drones. Buscamos garantizar el cumplimiento de la normativa vigente, prevenir incidentes y accidentes, proteger a las personas y bienes, y promover una cultura de seguridad en toda la organización.




% Capítulo 2: Politica y Objetivos
\newgeometry{top=4.2cm, bottom=2.5cm, left=2cm, right=2cm, headheight=60pt, headsep=1.5cm}
\chapter{Política y Objetivos}
\label{chap:Política y Objetivos}

\section{Politica y Compromiso de la Dirección}
SKYFILMS reconoce que la seguridad operacional es un valor fundamental y prioritario en todas nuestras actividades. La organización se compromete a demostrar un liderazgo activo y visible en la promoción de una cultura de seguridad positiva en toda la organización. Este compromiso se manifiesta a través de los siguientes principios que serán revisados anualmente:

\begin{enumerate}
    \item \textbf{Mejora Continua:} SKYFILMS se compromete a mejorar continuamente el nivel de rendimiento en materia de seguridad operacional, buscando activamente oportunidades para fortalecer nuestros procesos y reducir los riesgos.
    \item \textbf{Cultura de Seguridad:} Fomentar y mantener una cultura de seguridad operacional positiva en SKYFILMS, donde cada empleado se sienta responsable y partícipe en la identificación y gestión de riesgos.
    \item \textbf{Cumplimiento Normativo:} Cumplir con todos los requisitos normativos aplicables, incluyendo el RAC 100 y 219, así como cualquier otra regulación que afecte nuestras operaciones.
    \item \textbf{Recursos Adecuados:} Proporcionar los recursos humanos, financieros y técnicos necesarios para garantizar la seguridad de nuestras operaciones y el correcto funcionamiento del SMS.
    \item \textbf{Responsabilidad de la Gestión:} Asegurar que la seguridad operacional sea una responsabilidad primordial de todos los administradores y supervisores en SKYFILMS.
    \item \textbf{Comprensión e Implementación:} Garantizar que esta política de seguridad operacional sea comprendida, implementada y mantenida en todos los niveles de la organización.
    \item \textbf{Comportamiento inaceptable:} SKYFILMS define claramente los tipos de comportamientos que son inaceptables en relación con nuestras operaciones. Estos incluyen, pero no se limitan a:
    
        \begin{enumerate}[label=(\roman*)]
            \item Violaciones deliberadas de las regulaciones, políticas y procedimientos de seguridad.
            \item Conductas negligentes que pongan en riesgo la seguridad de las operaciones.
            \item Falsificación de registros o informes relacionados con la seguridad.
            \item Interferencia con los procedimientos de control de seguridad.
            \item Operar bajo la influencia de alcohol o sustancias psicoactivas.
        \end{enumerate}
    
    SKYFILMS se reserva el derecho de aplicar medidas disciplinarias proporcionales a la gravedad o falla en caso de comportamientos inaceptables. Sin embargo, no se aplicarán medidas en aquellos casos en que un empleado realice un informe de seguridad de buena fe, incluso si dicho informe revela una infracción no intencional de los procedimientos o regulaciones, siempre y cuando la acción no constituya negligencia grave o una violación deliberada.

\end{enumerate}

\vspace{1cm}

\begin{center}

\includegraphics[width=4cm]{Imagenes/firma.jpg}

\vspace{-1.2cm}

\rule{7cm}{0.5pt}

\vspace{0.3cm}
\textbf{Andrés Felipe Pérez Sierra}\\
C.C. 1.108.206.142\\
Gerente General\\


\end{center}
\newpage

\subsection{Compromiso de Notificación y Protección}
La Política invita a todo el personal a notificar peligros, errores y eventos, sin temor a represalias. Los reportes son confidenciales y no se usan para sancionar, salvo las excepciones de cultura justa descritas más abajo y en el Anexo 5. La Política también reconoce que cumplir el RAC 100 y el RAC 219 es obligatorio y que el SMS y sus procedimientos son de cumplimiento obligatorio para todo el personal.

\subsection{Aptitud para el Trabajo: Alcohol, Drogas y Fatiga}
\begin{itemize}
    \item \textbf{Cero tolerancia} al consumo de alcohol y sustancias psicoactivas (RAC 120): nadie opera, observa ni hace mantenimiento bajo sus efectos ni durante las horas previas. Si hay duda, no se trabaja y se informa al Gerente SMS.
    \item \textbf{Medicamentos:} quien tome un medicamento que cause sueño o afecte la atención lo informa al Jefe de Pilotos y no ocupa funciones críticas ese día.
    \item \textbf{Fatiga:} descanso de al menos 8 horas entre jornadas, jornada de vuelo y de mantenimiento fijada por el Jefe de Pilotos y el Gerente SMS, y derecho de cualquier persona a negarse a trabajar si está agotada. Quien lo haga no recibe ninguna sanción.
    \item \textbf{Funciones críticas} (piloto al mando, observador y técnico que libera aeronaves): antes de cada jornada, se declara con el FRAT (Anexo 6) que se está en condiciones de trabajar.
\end{itemize}

\subsection{Comunicación y Revisión de la Política}
\begin{enumerate}
    \item \textbf{Publicación:} el Gerente SMS entrega la Política firmada a todo el personal, y la deja visible en la oficina y en el repositorio digital. A clientes, proveedores y aseguradora se les entrega cuando lo soliciten o al contratar.
    \item \textbf{Comprensión:} en la inducción y en cada reentrenamiento se explica la Política y el personal firma que la recibió y la entendió (F-015). Se verifica con preguntas en el examen del Anexo 12.
    \item \textbf{Participación visible de la dirección:} el Gerente General la presenta personalmente en la reunión anual y firma las actas del SRB.
    \item \textbf{Eficacia:} en el SRB se revisa si el personal conoce la Política (resultados de evaluaciones y reportes) y si el medio de comunicación sirvió; si no, se cambia.
    \item \textbf{Revisión:} el Gerente General revisa la Política una vez al año y cuando cambie el alcance, y la vuelve a firmar. Se guardan las evidencias cinco años.
\end{enumerate}

\section{Objetivos de Seguridad Operacional}
Los objetivos del Sistema de Gestión de Seguridad Operacional (SMS) son fundamentales para guiar a la organización y establecer metas claras que permitan medir la eficacia y el desempeño del sistema. Para SKYFILMS, estos propósitos representan el compromiso de la alta dirección con la mitigación proactiva de los riesgos en todas nuestras operaciones con aeronaves no tripuladas (UAS).\\

\textbf{Objetivo General:}\\
Establecer, mantener y mejorar continuamente un Sistema de Gestión de Seguridad Operacional robusto que permita la identificación temprana de peligros, la mitigación efectiva de riesgos y la garantía de operaciones seguras, eficientes y en estricto cumplimiento normativo, manteniendo siempre los riesgos en un nivel aceptable.\\

\textbf{Objetivos Específicos:}\\
Para materializar nuestro objetivo general y asegurar la mejora continua del sistema, la organización ha definido los siguientes objetivos específicos:

\begin{itemize}
    \item \textbf{Mejora Continua:} Garantizar la implementación y efectividad del SMS a través de auditorías internas y revisiones periódicas por la dirección.

    \item \textbf{Cultura de Seguridad:} Fortalecer la cultura de seguridad operacional en todos los niveles mediante capacitación continua, talleres especializados y la promoción activa del reporte voluntario sin carácter punitivo.

    \item \textbf{Cumplimiento Normativo:} Asegurar el cumplimiento irrestricto de los estándares regulatorios aplicables, en particular las disposiciones del RAC 100, RAC 219 y las circulares informativas vigentes de la autoridad aeronáutica.

    \item \textbf{Gestión de Recursos Adecuados:} Proveer y gestionar de manera eficiente y oportuna los recursos humanos, técnicos y financieros necesarios para el correcto funcionamiento del SMS.
\end{itemize}



Con el fin de facilitar el seguimiento, control y la medición del desempeño de estas metas, SKYFILMS ha vinculado estos objetivos a Indicadores de Rendimiento de Seguridad Operacional (SPI) específicos y medibles. El desglose detallado de las metas, periodicidad, fórmulas de medición y planes de acción asociados a cada objetivo se encuentra documentado en el \textbf{Anexo 1 - Objetivos de Seguridad Operacional KPI's}.


\subsection{Objetivos, SPI y Metas}
Cada objetivo específico se relaciona con uno o más indicadores y metas, que están en el Anexo 1 y el Anexo 9. Para cumplir con el Programa Estatal de Seguridad Operacional, los objetivos se alinean con los del Estado (prevenir accidentes, reportar los eventos obligatorios y cumplir el RAC 100 y el RAC 219). El cumplimiento de los SPI hace parte de los objetivos específicos de Skyfilms. Los SPI cubren las operaciones certificadas: VLOS, vuelo nocturno, zona urbana y la flota Mavic 4 Pro y Pavo20 Pro II.

\subsection{Comunicación y Revisión de los Objetivos}
El Gerente General aprueba los objetivos. El Gerente SMS los comunica con la Política (mismo procedimiento) y los presenta al SRB una vez al año con los resultados de los SPI. Si un SPI no se cumple, se corrige o el objetivo se revisa y se deja constancia en el acta del SRB.

 \section{Cultura Justa en el Sistema de Gestión}
Para SKYFILMS, la Cultura Justa constituye la piedra angular sobre la cual se cimienta el Sistema de Gestión de Seguridad Operacional (SMS). De acuerdo con los principios del Documento 9859 de la OACI y los requisitos del RAC 219, la organización define la Cultura Justa como un entorno operativo en el cual se alienta, protege y capacita al personal (pilotos a distancia, observadores de espacio aéreo, personal de mantenimiento y directivos) para notificar de manera transparente cualquier información esencial relacionada con la seguridad operacional, garantizando que no serán objeto de castigos, represalias o acciones punitivas por sus acciones, omisiones o decisiones, siempre que estas correspondan a su experiencia y entrenamiento.\\

Este enfoque reconoce explícitamente que el error humano es una variable inevitable dentro de los sistemas tecnológicos y aeronáuticos complejos, incluidas las misiones comerciales de filmación y operaciones especiales con drones (UAS). Por consiguiente, la alta dirección de SKYFILMS se compromete formalmente a mantener un equilibrio ético y operativo entre dos ejes fundamentales:

\begin{itemize}
    \item \textbf{El Enfoque de Aprendizaje Organizacional:} Centrado en investigar y comprender qué falló en el sistema (defensas ausentes, fallos de equipos, factores ambientales o ambigüedad en los procedimientos) en lugar de buscar quién cometió el error, permitiendo extraer lecciones técnicas orientadas a la mejora continua.

    \item \textbf{La Exigencia de Responsabilidad (Accountability):} El establecimiento de fronteras claras e inequívocas que diferencien las fallas honestas del sistema de aquellos comportamientos que constituyen faltas disciplinarias o conductas inaceptables que comprometan la seguridad de las personas o bienes en superficie.
\end{itemize}

\subsection{Procedimiento de Cultura Justa}
 Ante la ocurrencia de cualquier evento de seguridad operacional, recepción de un reporte voluntario (VOR) u obligatorio (MOR), o hallazgo detectado en auditoría, el Gerente SMS y el Comité de Seguridad Operacional aplicarán formalmente el conducto regular analítico estructurado en las siguientes fases consecutivas:

 \definecolor{headerblue}{RGB}{180, 198, 231}
 
\renewcommand{\arraystretch}{1.4}
\setlength{\tabcolsep}{4pt}
 
\begin{table}[H]
\begin{adjustbox}{center}
\centering
\footnotesize
\begin{tabular}{|>{\centering\arraybackslash}m{1.1cm}|>{\centering\arraybackslash}m{2.9cm}|>{\RaggedRight\arraybackslash}m{5cm}|>{\centering\arraybackslash}m{3cm}|>{\centering\arraybackslash}m{2.3cm}|>{\centering\arraybackslash}m{3cm}|}
\hline
\rowcolor{gray!20}
\textbf{ÍTEM} & \textbf{ACTIVIDAD} & \textbf{DESCRIPCIÓN} & \textbf{RESPONSABLE} & \textbf{REGISTRO} & \textbf{FRECUENCIA} \\
\hline
 
1 &
Recepción del reporte &
Canalizar reportes sobre incidentes o situaciones inseguras, garantizando confidencialidad y libre expresión. &
Gerente SMS &
Formato de reporte &
Cada vez que se genere un reporte \\
\hline
 
2 &
Análisis preliminar &
Realizar un análisis inicial para determinar la necesidad de una investigación formal, según la naturaleza del evento. &
Gerente SMS &
Memo interno &
Según el caso \\
\hline
 
3 &
Investigación técnica &
Ejecutar una investigación objetiva que identifique causas y factores contribuyentes, enfocándose en el aprendizaje y no en la sanción. &
Gerente SMS &
Informe técnico &
Cada evento reportado \\
\hline
 
4 &
Clasificación de conducta &
Determinar la naturaleza del comportamiento (error, negligencia, violación) con base en criterios establecidos de cultura justa. &
Gerente SMS &
Matriz de clasificación &
Trimestral \\
\hline
 
5 &
Comité de revisión &
Evaluar los hallazgos y definir las acciones correspondientes: educativas, correctivas o disciplinarias si aplica. &
SAG -- Safety Action Group &
Acta de comité &
Según requerimiento \\
\hline
 
6 &
Comunicación de resultados &
Informar al personal implicado sobre las conclusiones del análisis, respetando la confidencialidad y promoviendo el aprendizaje. &
Gerente SMS &
Informe de retroalimentación &
Cuando se requiera \\
\hline
 
7 &
Seguimiento y verificación &
Supervisar el cumplimiento de las acciones definidas, verificando su impacto y cierre efectivo. &
Gerente SMS &
Consolidado de acciones &
Cuando aplique \\
\hline
 
8 &
Formación y sensibilización &
Implementar actividades de capacitación en cultura justa y seguridad operacional, dirigidas a todo el personal involucrado. &
Jefe de pilotos y gerente de SMS &
Registro de capacitaciones &
Semestral o según necesidad \\
\hline
 
9 &
Retroalimentación al sistema &
Documentar los aprendizajes obtenidos y proponer mejoras al SMS, MO o demás manuales aplicables. &
SAG &
Informe de mejora continua &
Semestral \\
\hline
 
\end{tabular}
\end{adjustbox}
\end{table}



\subsection{Protección de la Información y Tiempos de Gestión}
La protección de la identidad del reportante, la desidentificación de la información, las excepciones para divulgarla y los plazos de acuse, análisis, cierre y retroalimentación para cada tipo de reporte están definidos en el Anexo 5. La retroalimentación al personal se hace de forma individual al reportante y de forma general en el boletín de seguridad del trimestre y en la reunión de seguridad.

\textbf{Conductas que no se protegen (excepciones):} violación deliberada de normas, negligencia grave, falsificación de registros y operar bajo alcohol o sustancias psicoactivas. En esos casos, la decisión sobre la acción disciplinaria la toma el Gerente General con el informe del Gerente SMS.

\section{Rendición de Cuentas y Responsabilidades}
En SKYFILMS, la seguridad operacional es una responsabilidad compartida por todos los miembros de la organización. Sin embargo, existen roles y responsabilidades específicas asignadas a cada cargo para garantizar la implementación efectiva del Sistema de Gestión de Seguridad Operacional (SMS) y el cumplimiento de la normativa vigente.\\

\subsection{Gerente General:} El Gerente General es la máxima autoridad en SKYFILMS en lo que respecta a la seguridad operacional. Tiene la responsabilidad final en la conducción de los asuntos de certificación o aprobación de la organización, de las operaciones autorizadas y de todos los asuntos de seguridad operacional, incluyendo el establecimiento y la promoción de las Políticas del SMS. Además, es la autoridad final sobre las cuestiones de recursos humanos, técnicos y financieros necesarios para el rendimiento eficaz y eficiente del sistema de gestión.\\

\subsection{Gerente SMS:} El Gerente SMS es asignado para ejecutar y documentar todas las tareas de implementación, ejecución y mantenimiento del SMS en nombre del Gerente General. Sus responsabilidades incluyen:

\begin{itemize}
    \item Asegurarse de que los procesos del SMS se establezcan, implementen, ejecuten y mantengan.
    \item Promover la conciencia y una cultura positiva de seguridad operacional.
    \item Servir de enlace con la Aerocivil en cuestiones relacionadas con la seguridad operacional.
    \item Intercambiar lecciones valiosas aprendidas con otros operadores de UAS.
    \item Gestionar las investigaciones internas de incidentes y accidentes.
    \item Asegurarse de que se estén identificando los peligros y gestionando los riesgos de la operación del explotador UAS.
    \item Mantener la documentación de seguridad operacional.
    \item Organizar la instrucción en seguridad operacional.
\end{itemize}

\subsubsection{Perfil de Competencias}
La persona asignada como Gerente SMS demostrar conocimientos o formación en áreas del sector aeronáutico y contar con experiencia operacional aeronáutica de, al menos, doce (12) meses.\\

\textbf{NOTA:} De acuerdo con el RAC 1 “Definiciones”, “actividad aeronáutica” significa: “Conjunto de tareas y operaciones, directa o indirectamente relacionadas con el empleo de aeronaves civiles.”

\subsubsection{Designación, Dedicación y Recursos}
\begin{itemize}
    \item \textbf{Nombramiento:} el Gerente General designa por escrito al Gerente SMS, con una carta que indica su nombre, cargo y funciones: [nombre y cédula del Gerente SMS por diligenciar]. La carta se archiva como registro del SMS.
    \item \textbf{Línea jerárquica:} reporta directamente al Gerente General.
    \item \textbf{Dedicación y recursos:} dispone del tiempo necesario para sus funciones (mínimo medio día por semana, y más en las épocas de auditoría o evento), acceso al repositorio digital, apoyo de los demás cargos y presupuesto para capacitación.
    \item \textbf{Competencia:} formación en SMS (curso inicial y el curso de personal clave del Anexo 12), conocimiento de las operaciones UAS y experiencia aeronáutica de al menos doce meses, como se dijo arriba.
    \item \textbf{Conflicto de intereses:} no es Gerente General ni Jefe de Pilotos y no aprueba sus propias acciones. Si el evento le involucra directamente, el Gerente General designa otra persona para investigarlo.
\end{itemize}

\subsubsection{Autoridad para Detener Operaciones}
El Gerente General, el Gerente SMS y el Jefe de Pilotos pueden detener o cancelar un vuelo, una misión o una operación cuando haya un riesgo no aceptable. Cualquier piloto u observador puede detener su propio vuelo. Nadie recibe sanción por detener una operación por razones de seguridad. La reanudación la autoriza el Jefe de Pilotos cuando el riesgo vuelva a ser aceptable y, si la causa es una zona roja, el Gerente General.

\subsection{Jefe de Pilotos:} El Jefe de Pilotos es el encargado de supervisar y asegurar que todos los procedimientos establecidos en el manual de operaciones se cumplan según se haya establecido. Sus responsabilidades son:

\begin{itemize}
    \item Gestionar las solicitudes de autorización de vuelos UAS ante la Aerocivil.
    \item Garantizar que las operaciones del operador UAS se efectúen de manera segura, de acuerdo con el RAC 100 y demás normas vigentes aplicables.
    \item Asegurar la adecuada asignación y utilización de cada UA y de cada piloto UAS, de acuerdo con las competencias requeridas, para el desarrollo de la operación autorizada al operador UAS.
    \item Mantener un registro documental de las competencias de cada piloto UAS designado.
    \item Monitorizar la competencia de cada piloto UAS designado.
    \item Proponer las directrices y procedimientos estándar para las operaciones autorizadas a SKYFILMS, según corresponda, socializando, monitorizando y supervisando su cumplimiento por parte de cada piloto UAS designado y demás personal involucrado en las operaciones.
\end{itemize}

\subsection{Pilotos y Observadores:} Los pilotos y observadores son responsables de la operación directa de los UAS. Sus responsabilidades incluyen:

\begin{itemize}
    \item Realizar vuelos de manera segura y cumpliendo con la normativa.
    \item Mantener una comunicación efectiva con el control de tráfico aéreo y otros miembros del equipo.
    \item Reportar cualquier anomalía o condición insegura.
\end{itemize}

La asignación clara de responsabilidades y la rendición de cuentas son fundamentales para el éxito del SMS de SKYFILMS. Cada miembro del equipo es responsable de cumplir con sus obligaciones y de contribuir a la seguridad de las operaciones.

\section{Estructura Organizacional y Personal Clave}
SKYFILMS cuenta con una estructura organizacional jerárquica y funcionalmente definida para garantizar la eficiencia y seguridad de sus operaciones. Esta estructura facilita la comunicación, la toma de decisiones y la rendición de cuentas en todos los niveles. Además, dos grupos clave, el Grupo de Acción de Seguridad Operacional (SAG) y la Junta de Revisión de Seguridad Operacional (SRB), complementan esta estructura para una gestión integral de la seguridad.\\

En la cúspide se encuentra el \textbf{Gerente General}, quien tiene la autoridad final en todas las cuestiones de seguridad y es responsable del SMS. El \textbf{Gerente SMS} reporta directamente a él y se encarga de la implementación y mantenimiento del SMS.\\

El \textbf{Jefe de Pilotos} supervisa las operaciones de vuelo, asegurando el cumplimiento de procedimientos y normativas. Los \textbf{Pilotos} son responsables de la operación segura de los UAS y reportan al Jefe de Pilotos.\\

Para la gestión de riesgos y la mejora continua, SKYFILMS cuenta con dos grupos especializados:

\begin{itemize}
    \item La \textbf{Junta de Revisión de Seguridad Operacional (SRB)}, liderada por el Gerente General y conformada por el Gerente SMS y el Jefe de Pilotos, se enfoca en la gestión estratégica de la seguridad. Sus funciones incluyen la revisión del desempeño del SMS, la evaluación de riesgos de alto nivel, la asignación de recursos y la definición de directrices estratégicas para el SAG. Esta junta se reunirá con una periodicidad \textbf{anual} o cada que se identifique un peligro nuevo con un nivel de riesgo en la región \textbf{NO Tolerable} o \textbf{Tolerable} y cada reunión deberá tener un acta de realización emitida el \textbf{Gerente General}.
    \item El \textbf{Grupo de Acción de Seguridad Operacional (SAG)}, liderado por el Gerente SMS y conformado por los responsables de las áreas funcionales que influyen en la operación final, se encarga de la gestión táctica de la seguridad. Esto incluye la identificación de peligros, la evaluación y mitigación de riesgos, y la supervisión de la eficacia de las medidas de seguridad. Este grupo se reunirá cada que la Junta de Revisión de Seguridad Operacional lo solicite.
\end{itemize}

Esta estructura organizacional, que incluye los grupos SAG y SRB, asegura una clara definición de roles, facilita la comunicación y coordinación, y permite una gestión integral de la seguridad operacional en SKYFILMS. La tabla a continuación ilustra la jerarquía y las relaciones entre los diferentes cargos y grupos.

\begin{table}[H]
\centering
\renewcommand{\arraystretch}{1.35}
\begin{tabularx}{\textwidth}{|>{\RaggedRight\arraybackslash}p{4.4cm}|>{\RaggedRight\arraybackslash}X|>{\RaggedRight\arraybackslash}p{4.4cm}|}
\hline
\rowcolor[gray]{0.9}\textbf{Instancia} & \textbf{Función} & \textbf{Se relaciona con} \\ \hline
\rowcolor[HTML]{F4D6D0}Gerente General (Ejecutivo Responsable) & Máxima autoridad; toma las decisiones de seguridad operacional. & Preside el SRB \\ \hline
\rowcolor[HTML]{FFF2A8}Junta de Revisión de Seguridad Operacional (SRB) & Revisa el desempeño del SMS y decide sobre recursos y riesgos altos. & Gerente General; Gerente SMS; Jefe de Pilotos; SAG \\ \hline
\rowcolor[HTML]{FFF2A8}Gerente SMS (coordinador) & Administra el SMS y coordina las acciones. & SRB y SAG \\ \hline
Jefe de Pilotos & Supervisa la operación y aplica las acciones en vuelo. & SRB y SAG \\ \hline
\rowcolor[HTML]{FFF2A8}Grupo de Acción de Seguridad Operacional (SAG) & Ejecuta y da seguimiento a las acciones. & SRB; Gerente SMS; Jefe de Pilotos \\ \hline
\end{tabularx}
\caption{Estructura del SMS de Skyfilms.}
\end{table}


\subsection{Conflicto de Intereses}
En SKYFILMS como medida para evitar el conflicto de intereses se define que, el \textbf{Gerente SMS} no podrá ejercer en ningún momento el cargo de \textbf{Gerente General} ni \textbf{Jefe de pilotos}.


\subsection{Grupos de Seguridad: SRB y SAG}
\begin{longtable}{|>{\RaggedRight\arraybackslash}p{2.2cm}|>{\RaggedRight\arraybackslash}p{3.2cm}|>{\RaggedRight\arraybackslash}p{3.2cm}|>{\RaggedRight\arraybackslash}p{3.0cm}|>{\RaggedRight\arraybackslash}p{3.4cm}|}
\hline
\rowcolor[gray]{0.9}\textbf{Grupo} & \textbf{Integrantes} & \textbf{Frecuencia} & \textbf{Quórum} & \textbf{Qué hace y cómo se informa} \\ \hline
\endhead
SRB & Gerente General (preside), Gerente SMS, Jefe de Pilotos y Técnico UAS & Una vez cada seis meses y de inmediato ante un accidente, incidente grave o riesgo rojo & Gerente General y Gerente SMS & Revisa SPI, riesgos altos, recursos y auditorías. Se deja acta firmada y el Gerente SMS comunica las decisiones a todo el personal en cinco días hábiles \\ \hline
SAG & Gerente SMS (lidera), Jefe de Pilotos, Técnico UAS y un piloto & Cada trimestre y cuando haya un evento & Gerente SMS y dos integrantes más & Analiza reportes, riesgos y acciones; da seguimiento y escala al SRB lo que exceda su autoridad \\ \hline
\end{longtable}

\subsection{Niveles de Autoridad para Aceptar Riesgos}
\begin{longtable}{|>{\RaggedRight\arraybackslash}p{3cm}|>{\RaggedRight\arraybackslash}p{4.3cm}|>{\RaggedRight\arraybackslash}p{4.4cm}|>{\RaggedRight\arraybackslash}p{3.7cm}|}
\hline
\rowcolor[gray]{0.9}\textbf{Región} & \textbf{Quién decide} & \textbf{Qué se hace y en qué plazo} & \textbf{Revisión} \\ \hline
\endhead
Roja (no tolerable) & Gerente General & La operación se detiene de inmediato. Plan de acción en 24 horas. No se reanuda hasta que el riesgo residual sea amarillo o verde & Seguimiento semanal hasta el cierre \\ \hline
Amarilla con índice 4C, 4D, 4E, 5D, 5E o 3B & Gerente SMS, con informe al Gerente General en el SRB & Plan de acción en 5 días hábiles; controles puestos en 15 días hábiles & Cada mes \\ \hline
Amarilla (otros índices) & Gerente SMS, con visto bueno del Jefe de Pilotos & Plan de acción en 10 días hábiles; controles en 30 días hábiles & Cada trimestre \\ \hline
Verde (aceptable) & Gerente SMS & Se acepta y se vigila & Cada año \\ \hline
\end{longtable}

\section{Coordinación de Respuesta Ante Emergencias}
SKYFILMS, reconoce la importancia de contar con un plan de respuesta a emergencias (PRE) sólido y efectivo. Este PRE, en línea con las regulaciones colombianas, tiene como objetivo establecer los procedimientos y acciones necesarias para responder de manera oportuna y eficiente ante cualquier eventualidad que pueda afectar la seguridad de las operaciones, el personal, los equipos o el medio ambiente.

\subsection{Alcance}
Este PRE abarca todas las operaciones aéreas realizadas por SKYFILMS con sus UAS, así como las actividades desarrolladas en sus instalaciones y el bienestar de su personal. Esta sección debe ser conocida y consultada por todo el personal en caso de emergencia.

\subsection{Objetivos}

\begin{itemize}
    \item Proteger la vida e integridad de las personas involucradas en las operaciones de SKYFILMS.
    \item Minimizar los daños a los equipos, las instalaciones y el medio ambiente.
    \item Garantizar la continuidad de las operaciones en la medida de lo posible.
    \item Facilitar la comunicación y coordinación con las autoridades competentes y los servicios de emergencia.
    \item Establecer los procedimientos para la investigación de incidentes y la implementación de medidas correctivas.
\end{itemize}

\subsection{Marco Legal}
Este PRE se fundamenta en las siguientes normativas:

\begin{itemize}
    \item Reglamentos Aeronáuticos de Colombia (RAC) 100
    \item Reglamentos Aeronáuticos de Colombia (RAC) 219
\end{itemize}

\subsection{Estructura Organizacional de Respuesta a Emergencias}
La estructura de respuesta a emergencias de SKYFILMS se basa en los siguientes roles y responsabilidades:

\begin{itemize}
    \item \textbf{Gerente General:}
        \begin{itemize}
            \item Autoridad máxima en la toma de decisiones durante una emergencia.
            \item Responsable de la activación y ejecución del PRE.
            \item Enlace con las autoridades y los medios de comunicación.
        \end{itemize}
    \item \textbf{Gerente SMS:}
        \begin{itemize}
            \item Responsable de la coordinación general de la respuesta a la emergencia.
            \item Encargado de la comunicación interna y externa.
            \item Responsable de la investigación de incidentes y la implementación de medidas correctivas.
        \end{itemize}
    \item \textbf{Jefe de Pilotos:}
        \begin{itemize}
            \item Responsable de la seguridad de las operaciones de vuelo.
            \item Encargado de la coordinación de los pilotos y observadores durante una emergencia.
        \end{itemize}
    \item \textbf{Pilotos y Observadores:}
        \begin{itemize}
            \item Responsables de seguir los procedimientos de emergencia establecidos.
            \item Encargados de la comunicación con el Jefe de Pilotos y el Gerente SMS.
        \end{itemize}
\end{itemize}

\subsection{Procedimientos de Respuesta a Emergencias}

\subsubsection{Detección y Alerta}
\begin{itemize}
    \item Cualquier persona que presencie o tenga conocimiento de una emergencia deberá comunicarlo de inmediato al \textbf{Jefe de Pilotos} o, en su ausencia, al \textbf{Gerente SMS.}
    \item La comunicación se realizará utilizando los medios disponibles más rápidos (teléfono, radio, correo electrónico, etc.).
    \item Se proporcionará la siguiente información:
        \begin{itemize}
            \item Naturaleza de la emergencia
            \item Ubicación exacta
            \item Número de personas afectadas (si aplica)
            \item Gravedad de la situación
            \item Información de contacto del reportante
        \end{itemize}
\end{itemize}

\subsubsection{Evaluación Inicial}
\begin{itemize}
    \item El \textbf{Jefe de Pilotos} o el \textbf{Gerente SMS} evaluarán la información recibida y determinarán la magnitud de la emergencia.
    \item Se considerarán factores como la gravedad de la situación, el potencial de daños y la necesidad de activar el PRE.
\end{itemize}

\subsubsection{Activación del PRE}
\begin{itemize}
    \item Si la evaluación inicial indica que la situación lo amerita, el \textbf{Jefe de Pilotos} o el \textbf{Gerente SMS} informará al \textbf{Gerente General} y declarará la activación del PRE.
    \item Se informará de inmediato a todo el personal relevante y se iniciarán los procedimientos específicos correspondientes a la emergencia en cuestión.
\end{itemize}

\subsubsection{Notificación a Autoridades}
\begin{itemize}
    \item En caso de ser necesario, se notificará a las autoridades competentes, como la Aerocivil, la policía, los bomberos, los servicios de emergencia médica, etc.
    
    \item La notificación se realizará utilizando los números de contacto designados en el directorio telefónico del PRE \textbf{(Anexo 2 - Contactos de emergencia)}.
    \item Se proporcionará información detallada sobre la emergencia, incluyendo la ubicación, la naturaleza, la gravedad y el número de personas afectadas.
\end{itemize}

\subsubsection{Procedimiento de Comunicación Directa con Aerocivil}

Durante la ocurrencia de una emergencia operacional, incidente grave o accidente que involucre las aeronaves no tripuladas (UAS) de SKYFILMS, es imperativo establecer un flujo de comunicación oficial, ordenado y oportuno con la Aerocivil (Aerocivil). Este procedimiento garantiza la correcta articulación institucional y previene la divulgación de información inexacta o no confirmada.\\

\textbf{Personal Autorizado para la Comunicación}

Para mantener una línea de mando unificada y proteger la integridad de la información durante la crisis, se restringe la vocería oficial y el contacto directo con la autoridad aeronáutica exclusivamente a los siguientes cargos:

\begin{itemize}
    \item \textbf{Gerente General:} Como Ejecutivo Responsable ante la autoridad, encargado de las notificaciones corporativas y legales.

    \item \textbf{Gerente SMS:} Como enlace técnico principal para el reporte de la emergencia y el flujo de información de seguridad.

    \item \textbf{Jefe de Pilotos:} Como apoyo técnico para detallar aspectos operacionales específicos del evento ante los inspectores de la autoridad.
\end{itemize}

Ningún otro miembro del personal operativo, técnico o administrativo de SKYFILMS está autorizado para emitir declaraciones, entregar reportes preliminares o establecer contacto oficial con la Aerocivil durante una emergencia, debiendo escalar cualquier requerimiento a los cargos mencionados.\\

\textbf{Fines Específicos de la Comunicación}\\

El contacto de emergencia con la Aerocivil se activará de forma inmediata y prioritaria, ciñéndose estrictamente a los siguientes propósitos operativos y regulatorios:

\begin{itemize}
    \item \textbf{Coordinación Operativa:} Notificación y articulación técnica con el Grupo de Drones (UAS) de la Aerocivil respecto a la condición de la aeronave, el espacio aéreo afectado y las medidas iniciales de contención adoptadas por la empresa.

    \item \textbf{Búsqueda y Rescate (SAR):} Solicitud de apoyo, soporte y coordinación con los servicios de Búsqueda y Salvamento en caso de extravío de la aeronave en zonas críticas o afectación a personas en superficie.

    \item \textbf{Investigación de Accidentes:} Notificación inicial para la activación de los protocolos de preservación de evidencia y articulación con el grupo de Investigación de Accidentes e Incidentes de Aviación de la Aerocivil.
\end{itemize}


\textbf{Canales de Comunicación}
Las comunicaciones se realizarán única y exclusivamente a través de los canales de contacto de emergencia oficiales, directorios y líneas telefónicas actualizadas y publicadas en la página web oficial de la Aerocivil. Para garantizar una respuesta ágil, el Gerente SMS será responsable de verificar semestralmente la vigencia de estos números y mantenerlos actualizados en el \textbf{Anexo 2 - Contactos de Emergencia} del presente manual.\\

Para llevar a cabo la comunicación con Aerocivil se deberán ejecutar los siguientes pasos: 



\begin{table}[H]
\begin{adjustbox}{center}
\centering
\footnotesize
\begin{tabular}{|>{\centering\arraybackslash}m{0.06\linewidth}|m{0.55\linewidth}|>{\centering\arraybackslash}m{0.25\linewidth}|}
\hline
\rowcolor{gray!20}
\textbf{PASO} & \textbf{ACTIVIDAD} & \textbf{RESPONSABLE} \\
\hline

1 &
Identificar la situación de emergencia. &
Piloto UAS, Observador UAS \\
\hline

2 &
Reconocer la necesidad de contactar a la Aerocivil debido a emergencias, accidentes o incidentes. &
Jefe de Pilotos, Gerente SMS \\
\hline

3 &
Evaluar el canal de comunicación mas eficiente, considerando la urgencia para la toma de decisiones, el cual puede ser por medio telefónico o a través de correo electrónico. &
Jefe de Pilotos, Gerente SMS \\
\hline

4 &
Establecimiento de contacto inicial, Informando la naturaleza del mensaje y proporcionando detalles relevantes, como la ubicación y el número de identificación del UA. &
Gerente \\
\hline

5 &
Confirmación y registro de la comunicación, incluyendo la hora, la naturaleza del mensaje y cualquier acción tomada como resultado de la comunicación. &
Jefe de Pilotos, Gerente SMS \\
\hline

\end{tabular}
\end{adjustbox}
\caption{Procedimiento de Comunicación con Aerocivil}
\label{tab:responsabilidades_documentales}
\end{table}

\subsubsection{Respuesta Inicial}
\begin{itemize}
    \item Se implementarán las acciones inmediatas necesarias para controlar la situación, minimizar los daños y proteger a las personas y los bienes. Esto puede incluir:
    \begin{itemize}
        \item Suspensión de operaciones de vuelo.
        \item Aislamiento del área afectada.
        \item Evacuación de instalaciones.
        \item Prestación de primeros auxilios.
        \item Activación de equipos de respuesta a emergencias (si aplica).
    \end{itemize}
\end{itemize}

\subsubsection{Investigación}
\begin{itemize}
    \item Se llevará a cabo una investigación exhaustiva liderada por el Gerente SMS para determinar las causas de la emergencia y se documentarán los hallazgos.
    \item Se tomarán medidas correctivas y preventivas para evitar futuras ocurrencias y se informara por medio de boletines las lecciones aprendidas y/o retroalimentación.
\end{itemize}

\subsubsection{Comunicación con Medios}
\begin{itemize}
    \item En caso de que la emergencia tenga repercusión mediática, los únicos autorizados para comunicarse con los medios de comunicación son el \textbf{Gerente General, Gerente SMS y Jefe de Pilotos}.
    \item Se proporcionará información precisa, oportuna y transparente, evitando especulaciones o declaraciones que puedan perjudicar la imagen de SKYFILMS.
\end{itemize}

%\subsubsection{Apoyo a Colaborador y Familias}
%\begin{itemize}
  %  \item Se brindará apoyo a los colaboradores y sus familias en caso de que sean afectados por la emergencia.
    %\item Esto puede incluir asistencia psicológica, legal o económica, según sea necesario.
%\end{itemize}

\subsubsection{Revisión y Mejora}
\begin{itemize}
    \item El PRE será revisado y actualizado periódicamente, y siempre después de una emergencia, para asegurar su eficacia y adaptabilidad.
    \item Se realizarán simulacros y ejercicios de emergencia de manera regular para capacitar al personal y poner a prueba la efectividad del plan.
\end{itemize}

\subsection{Procedimiento Específico para Emergencias con Baterías (UAS)}
Considerando que las aeronaves no tripuladas operadas por SKYFILMS utilizan baterías recargables de alta densidad energética (principalmente baterías de litio), se establece el siguiente procedimiento específico para la prevención, control y respuesta ante emergencias asociadas a baterías, tales como sobrecalentamiento, incendio, explosión, fuga térmica o daño estructural.

\subsubsection{Identificación del Riesgo}
Se consideran situaciones de emergencia asociadas a baterías, entre otras:

\begin{itemize}
    \item Inflado, deformación o daño físico visible en la batería.

    \item Sobrecalentamiento anormal durante carga, descarga o posterior al vuelo.

    \item Emisión de humo, olor químico fuerte o presencia de chispa.

    \item Incendio o combustión espontánea.

    \item Impacto severo durante la operación.

    \item Caída del UAS con posible afectación del compartimiento de baterías.
\end{itemize}

\subsubsection{Procedimiento de Actuación en Caso de Emergencia}
En caso de identificarse una situación anormal con la batería:

\begin{enumerate}
    \item \textbf{Durante la operación en vuelo:}
        \begin{enumerate}
            \item Activar procedimiento de aterrizaje inmediato en zona segura.

            \item Priorizar área libre de personas, vehículos e infraestructura crítica.

            \item En caso de pérdida de control asociada a falla energética, activar procedimiento de contingencia establecido.
        \end{enumerate}
    
    \item \textbf{En tierra (durante carga o manipulación):}
        \begin{enumerate}
            \item Suspender inmediatamente el proceso de carga.

            \item Aislar la batería en un área despejada y ventilada.

            \item No manipular directamente la batería si presenta incendio activo.

            \item Utilizar extintor tipo ABC o agente adecuado para fuegos clase B/C.

            \item Notificar inmediatamente al responsable de seguridad operacional.

            \item Si el evento escala, activar el Plan de Respuesta a Emergencias y coordinar con los servicios de emergencia externos (Bomberos, autoridad local).
        \end{enumerate}
\end{enumerate}

\subsubsection{Disposición Final de Baterías Dañadas}
Las baterías que hayan presentado:

\begin{itemize}
    \item Incendio.

    \item Inflado severo.

    \item Daño estructural significativo.

    \item Falla eléctrica recurrente.
\end{itemize}
Deberán ser retiradas permanentemente del servicio y dispuestas conforme a la normativa ambiental vigente para residuos peligrosos.\\

\textbf{NOTA:} Este procedimiento forma parte integral del Plan de Respuesta Ante Emergencias de SKYFILMS y es de cumplimiento obligatorio para todo el personal operativo y técnico involucrado en la operación de UAS.

\subsection{Procedimientos Específicos: Pérdida de Control y Pérdida de Enlace}
Los pasos detallados están en el Manual de Operaciones; aquí se resumen:
\begin{enumerate}
    \item El piloto avisa en voz alta al observador y al equipo y mantiene la calma.
    \item Intenta recuperar el enlace o el control; si hay retorno automático, lo deja actuar y no vuela hacia personas.
    \item El observador sigue la aeronave con la vista y anota el último punto y la hora.
    \item Si cae o desaparece, se asegura el lugar, se pone a salvo a las personas y se llama al 123 si hay riesgo.
    \item El Jefe de Pilotos informa al Gerente SMS y este decide el reporte a la Aerocivil (Anexo 4).
\end{enumerate}

\subsection{Pruebas y Simulacros del Plan de Respuesta}
\begin{itemize}
    \item \textbf{Frecuencia:} un simulacro cada seis meses, que se alterna entre escenarios de batería en llamas, pérdida de control y lesiones a una persona.
    \item \textbf{Cómo se hace:} el Gerente SMS elige el escenario y lo avisa al personal con anticipación; el Jefe de Pilotos coordina; el piloto y el observador actúan según el plan.
    \item \textbf{Evaluación:} el Gerente SMS mide tiempos, comunicación y uso del equipo, y escribe un informe de una página con los resultados.
    \item \textbf{Mejora:} las debilidades pasan al Plan de Mejora Continua (Anexo 11) y se corrige el plan, si hace falta, en 30 días.
    \item \textbf{Distribución y coordinación:} el plan se entrega a todo el personal con el Manual. Cuando operemos con un cliente o una entidad externa (por ejemplo, un lugar con seguridad propia), el Gerente SMS acuerda antes con ella quién llama a quién.
    \item \textbf{Otras regiones:} para cada ciudad o región nueva, antes de operar el Gerente SMS agrega los contactos al Anexo 2 y confirma el hospital y los bomberos más cercanos.
\end{itemize}

\subsection{Información de Contactos}
SKYFILMS mantendrá actualizada la información de los contactos de emergencias para las ciudades principales de Colombia en el \textbf{Anexo 2 - Contactos de Emergencia} y para las operaciones fuera de las ciudades principales de Colombia, previo a la operación, el \textbf{Gerente SMS} sera responsable de contactar las autoridades del sitio.


\section{Documentación SMS}
El acceso a la documentación del Sistema de Gestión de Seguridad Operacional (SMS) de SKYFILMS es fundamental para garantizar la trazabilidad, la transparencia y la eficacia de nuestros procesos de seguridad. Toda la documentación que aporta al sistema de seguridad operacional lleva una trazabilidad y se almacena de manera digital, asegurando que la información esté actualizada y sea accesible.

\subsection{Tipos de Documentos}
Cada documento del SMS de SKYFILMS se identifica con un código único según su naturaleza, el cual incluye la versión del documento. Estos documentos se almacenan en una carpeta principal designada para el SMS, que contiene subcarpetas para cada manual principal: Manual SMS, Manual de Operaciones y Manual de Mantenimiento. Dentro de cada subcarpeta, se almacena la trazabilidad de los documentos específicos que cada manual requiere para su cumplimiento.

\subsubsection{Categorización y Nomenclatura Documental}
Para garantizar la protección de la información sobre seguridad operacional y fuentes conexas, todos los documentos del SMS de SKYFILMS se clasifican en tres niveles de acceso, identificados mediante una estructura alfanumérica estándar:\\

\begin{table}[H]
\begin{adjustbox}{center}
\begin{tabular}{|>{\centering\arraybackslash}m{0.3\linewidth}|>{\centering\arraybackslash}m{0.45\linewidth}|}
\hline
\rowcolor{gray!20}
\textbf{Nomenclatura} & \textbf{Categoría} \\
\hline

PUB & Publico \\
\hline

INT & Interno \\
\hline

CON & Confidencial \\
\hline

\end{tabular}
\end{adjustbox}
\caption{Nomenclatura y Categorización }
\label{tab:responsabilidades_documentales}
\end{table}

\begin{table}[H]
\begin{adjustbox}{center}
\begin{tabular}{|>{\centering\arraybackslash}m{0.3\linewidth}|>{\centering\arraybackslash}m{0.45\linewidth}|}
\hline
\rowcolor{gray!20}
\textbf{Rol} & \textbf{Nivel de Acceso} \\
\hline

Gerente General & CON \\
\hline

Gerente SMS & CON \\
\hline

Jefe de Pilotos & CON \\
\hline

Pilotos UAS & INT \\
\hline
Técnico UAS & INT \\
\hline

Clientes / Aseguradora / Proveedores & PUB \\
\hline

\end{tabular}
\end{adjustbox}
\caption{Nivel de Acceso por Rol}
\label{tab:responsabilidades_documentales}
\end{table}

Los niveles de acceso (PUB, INT y CON) dicen quién puede \textit{ver} cada documento. Quién lo elabora, revisa, aprueba o custodia (E, R, A, C, S) está en la matriz del Anexo 3. Son dos matrices distintas y no se mezclan.

\subsection{Control de la Documentación}
La gestión de la documentación y los registros operacionales en SKYFILMS garantiza que el personal tenga acceso a la información actualizada, que las operaciones se realicen bajo los estándares vigentes y que exista trazabilidad para la gestión de riesgos y auditorías de la Aerocivil, para ello se crea el \textbf{Anexo 3 - Procedimiento de Control de Documentos y Registros}.

\subsubsection{Directrices Generales de Gestión Documental}
Todos los documentos de SKYFILMS deben regirse por las siguientes directrices obligatorias:

\begin{itemize}
    \item \textbf{Identificación Estándar:} Todo documento debe incluir en su encabezado: Código de identificación, título del documento, logotipo de la empresa, fecha de emisión, número de versión y paginación (Ej. Página X de Y).
    
    \item \textbf{Aprobación y Autoría:} Cada documento o revisión debe tener claramente definido el responsable de elaboración (Autor), revisión y aprobación (Ej. elaborado por Gerente SMS, aprobado por Ejecutivo Responsable).
    
    \item \textbf{Control de Cambios:} Deberá existir un cuadro de control de revisiones y enmiendas (Listado de Páginas Efectivas), registrando la fecha de modificación y la justificación del cambio.
    
    \item \textbf{Trazabilidad y Retención:} Los registros operacionales (como bitácoras de vuelo y registros de mantenimiento) deben conservarse disponibles durante un periodo mínimo de 2 años en estado activo, y 3 años adicionales en archivo inactivo (archivo muerto).

    \item \textbf{Preservación de Evidencia:} En caso de incidente o accidente, el registro de los logs de vuelo, audios y videos de la estación de control deben preservarse y bloquearse de inmediato para proteger la información ante la Dirección Técnica de Investigación de Accidentes (DIACC).
\end{itemize}

\subsubsection{Lista Maestra de Documentos SMS}
El SDCPS (Sistema de recopilación y procesamiento de datos sobre seguridad operacional) de SKYFILMS se compone de los siguientes manuales, procedimientos y registros:

\begin{table}[H]
\begin{adjustbox}{center}
\begin{tabular}{|>{\centering\arraybackslash}m{0.08\linewidth}|m{0.70\linewidth}|>{\centering\arraybackslash}m{0.1\linewidth}|}
\hline
\rowcolor{gray!20}
\textbf{No} & \textbf{Documento} & \textbf{Nivel de Acceso} \\
\hline

1 & Manual del Sistema de Gestión de Seguridad Operacional & INT \\
\hline

2 & Manual de Operaciones UAS & INT  \\
\hline

3 & Plan de Respuesta ante Emergencias (PRE / ERP). & INT  \\
\hline

4 & Procedimientos Operacionales Normales y Anormales de aviación no tripulada. & INT \\
\hline

5 & Matriz de Identificación de Peligros y Evaluación de Riesgos & CON \\
\hline

6 & Actas de Comités de Seguridad Operacional (SRB estratégico y SAG táctico). & CON  \\
\hline

7 & Actas de reuniones de pilotos y socialización de boletines de seguridad. & INT \\
\hline

8 & Formularios de Reporte Obligatorio (MOR) y Reporte Voluntario (VOR). & INT \\
\hline

9 & Autorizaciones de vuelo UAS (emitidas por la Aerocivil) y planes de vuelo internos. & INT \\
\hline

10 & Libros de Vuelo por cada UA (Logbooks individuales por dron). & INT \\
\hline

11 & Bitácoras de vuelo por cada Piloto UAS. & INT \\
\hline

12 & Manual de Control de Mantenimiento (MCM) y formatos de registro de mantenimiento técnico. & INT \\
\hline

13 & Manuales técnicos del fabricante del UAS y de los Equipos Tecnológicos Asociados (ETA). & INT \\
\hline

14 & Registros de Instrucción, Capacitación y Evaluación de Competencias del personal en SMS. & INT \\
\hline

15 & Indicadores de Rendimiento de Seguridad Operacional (SPIs, KPIs y Metas SPT). & CON \\
\hline

16 & Informes de Auditorías Internas y Externas del SMS. & CON \\
\hline

\end{tabular}
\end{adjustbox}
\caption{Lista Maestra de Documentos SMS}
\label{tab:responsabilidades_documentales}
\end{table}

\subsubsection{Asignación de Responsabilidades Documentales}
Para evitar cruces de funciones (conflicto de intereses) exigidos por el RAC 100 y RAC 219, la responsabilidad de la administración, diligenciamiento y custodia de los documentos se divide de la siguiente manera:

\begin{table}[H]
\begin{adjustbox}{center}
\centering
\renewcommand{\arraystretch}{1.5}
\begin{tabular}{|m{0.3\textwidth}|m{0.7\textwidth}|}
\hline
\rowcolor{gray!20}
\textbf{Rol Responsable} & \textbf{Documentos y Registros a su cargo} \\ \hline
\textbf{Gerente SMS} & 
\textbullet~Manual del SMS (MSMS).\newline
\textbullet~Matriz de Peligros y Riesgos (SRM).\newline
\textbullet~Actas de reuniones de comités SMS (SRB / SAG).\newline
\textbullet~Formularios y bases de datos de Reportes VOR / MOR.\newline
\textbullet~Indicadores de Rendimiento (SPIs y KPIs) y análisis de tendencias.\newline
\textbullet~Registros de capacitación e instrucción SMS del personal.\newline
\textbullet~Informes de investigación de eventos y auditorías internas. \\ \hline
\textbf{Jefe de Pilotos UAS} & 
\textbullet~Manual de Operaciones (MO).\newline
\textbullet~Libros de vuelo individuales por cada aeronave no tripulada (UA).\newline
\textbullet~Bitácoras de control de horas de vuelo de los Pilotos UAS.\newline
\textbullet~Procedimientos operacionales, Instructivos y Listas de Chequeo (Checklists).\newline
\textbullet~Planificaciones y Autorizaciones de Vuelo UAS ante Aerocivil.\newline
\textbullet~Certificados de idoneidad, médicas y adiciones de su equipo de pilotos. \\ \hline
\end{tabular}
\end{adjustbox}
\caption{Asignación de Responsabilidades Documentales del SMS}
\label{tab:responsabilidades_documentales}
\end{table}






% Capítulo 3: Gestion de Riesgos de Seguridad Operacional
\newgeometry{top=4.2cm, bottom=2.5cm, left=2cm, right=2cm, headheight=60pt, headsep=1.5cm}
\chapter{Gestión de Riesgos de Seguridad Operacional}
\label{chap:Gestión de Riesgos de Seguridad Operacional}
Para SKYFILMS, la gestión de riesgos de seguridad operacional es un pilar fundamental para garantizar la seguridad de nuestras operaciones con UAS. Reconocemos que el entorno aeronáutico presenta peligros inherentes, y nuestra capacidad para identificarlos y gestionar los riesgos asociados es crucial para prevenir incidentes y accidentes. Este proceso se fundamenta en un enfoque sistemático y formal que involucra la identificación de peligros, la evaluación, mitigación y aceptación de los riesgos de seguridad operacional.\\

La efectividad de esta gestión reside en la participación activa y el compromiso de cada individuo dentro de nuestra organización, desde la alta dirección hasta cada piloto. Para ello, SKYFILMS ha establecido un proceso formal y continuo para la identificación de peligros, combinando métodos reactivos, proactivos y predictivos, asegurando una cobertura integral de los posibles riesgos.\\

\section{Identificación de Peligros}
En SKYFILMS, la identificación de peligros es el primer y crucial paso en nuestra gestión de la seguridad operacional, entendiendo un peligro como una condición u objeto que podría provocar un incidente o accidente o contribuir al mismo. Nos centramos en una amplia gama de factores que pueden impactar nuestras operaciones con UAS, garantizando así el cumplimiento de los principios del SMS. Para ello, se consideran los siguientes criterios principales:

\begin{enumerate}
    \item \textbf{Peligros al personal operativo:} Condiciones en la operación que afectan directamente la integridad física y mental de nuestro equipo.
    \item \textbf{Peligros a personas ajenas a la operación:} Potenciales amenazas para la seguridad de terceros en tierra o en el aire.
    \item \textbf{Peligros en tierra (obstáculos naturales y artificiales):} Identificación de estructuras, terrenos, flora o fauna que puedan interferir con la operación.
    \item \textbf{Peligros en el aire (por fallas en la operación):} Riesgos derivados de posibles mal funcionamiento del equipo, errores humanos o condiciones operacionales adversas durante el vuelo.
    \item \textbf{Peligros estratégicos (por la interacción en espacios aéreos de aeronaves tripuladas):} Peligros asociados a la coexistencia con aeronaves tripuladas en el espacio aéreo, especialmente en zonas de operación compartidas.
    \item \textbf{Peligros por cambios en la normativa vigente aplicable:} La necesidad de adaptar constantemente nuestros procedimientos a las modificaciones en el RAC 100 y RAC 219, y otras regulaciones relevantes.
    \item \textbf{Peligros por la gestión del cambio:} Evaluación de los riesgos introducidos por cualquier modificación en nuestros procesos, equipos o personal, antes de su implementación.
\end{enumerate}

Para mantener nuestro Sistema de Gestión de Seguridad Operacional actualizado y en constante búsqueda de posibles peligros que aumenten el riesgo de las operaciones, SKYFILMS utiliza los métodos de identificación expuestos a continuación.

\subsection{Método Reactivo}
Este método SKYFILMS lo aplica para identificar peligros y factores contribuyentes a partir de eventos que ya han ocurrido. Esto nos permite aprender de la experiencia y evitar la recurrencia de sucesos no deseados. Las principales herramientas reactivas de SKYFILMS son:\\

\textbf{a) Sistema de Notificaciones de Seguridad Operacional:} Es la principal fuente de información reactiva. Fomentamos activamente una cultura de notificación, garantizando un tratamiento confidencial, protección de datos y no punitivo para los reportes voluntarios, esta información es almacenada en archivo digital con el objetivo de tener trazabilidad y tener presente la condición o peligro identificado para su retroalimentación y correcta gestión. Este sistema se divide en:

\begin{itemize}
    \item \textbf{Notificaciones Obligatorias (MOR):} Conforme a la regulación, SKYFILMS notificará de forma obligatoria a la Aerocivil (Grupo de Investigación de Accidentes, DIACC), dentro de las 48 horas siguientes (RAC 100.240), cualquier accidente o incidente grave. Así mismo, se reportarán los eventos de seguridad operacional definidos en la normativa vigente \textbf{(Anexo 4 - Listado Reportes Obligatorios M.O.R)}. Toda notificación obligatoria amerita una investigación formal por parte de la Junta de Revisión de Seguridad Operacional SRB.
    \item \textbf{Notificaciones Voluntarias:} Se incentiva a todo el personal y a terceros a reportar cualquier condición, acto o situación que perciban como un peligro potencial, incluso si no resultó en un incidente. Esto es crucial para identificar condiciones latentes y tendencias antes de que escalen a eventos mayores, siguiendo la "Teoría del Iceberg". Para facilitar este proceso, se utilizará un formato estandarizado.
\end{itemize}

\textbf{NOTA 1:} En cualquier caso de notificación, ya sea obligatoria o voluntaria se realizara retroalimentación presencial o virtual a la persona que realizo la notificación, la \textbf{Junta de Revisión de Seguridad Operacional SRB} es la encargada de analizar dicha notificación y la retroalimentación estará a cargo del \textbf{Gerente SMS}.\\

\textbf{NOTA 2:} En el \textbf{“Anexo 5 - Guía de Uso Formato de Notificación Peligros de Seguridad Operacional”} se explica el método de uso del formato de notificaciones.\\

\textbf{b) Investigación de Sucesos de Seguridad Operacional:} No todas las notificaciones requerirán una investigación a fondo. El \textbf{Gerente SMS}, en conjunto con la \textbf{Junta de Revisión de Seguridad Operacional (SRB)}, clasificará los reportes según el nivel de riesgo potencial para determinar cuáles ameritan una investigación formal. El objetivo de toda investigación no es encontrar culpables, sino determinar las causas raíz (factores contribuyentes y organizacionales) para implementar defensas efectivas. Para ello, SKYFILMS podrá utilizar técnicas de análisis reconocidas como:

\begin{itemize}
    \item Los 5 Porqués.
    \item Diagrama de Ishikawa (Causa-Efecto o Espina de Pescado).
    \item Análisis Bow-Tie (Corbatín).
\end{itemize}

Los hallazgos de cada investigación se documentarán y se utilizarán para actualizar el registro de peligros y los procesos de gestión de riesgos.

\subsection{Métodos Proactivo y Predictivo}
Estos métodos se enfocan en la anticipación, buscando peligros mediante el análisis de las operaciones actuales (proactivo) y la proyección de problemas futuros a partir de datos históricos (predictivo). SKYFILMS los implementa para crear defensas robustas antes de que los peligros se materialicen de la siguiente manera:\\

\textbf{a) Inspecciones y Auditorías de Seguridad Operacional:}

\begin{itemize}
    \item \textbf{Auditorías Internas:} Se realizarán auditorías planificadas a los procesos del SMS y del Manual de Operaciones (MO) para verificar el cumplimiento y la efectividad de los procedimientos.
    \item \textbf{Inspecciones Pre-vuelo:} Antes de cada vuelo, el piloto al mando realizará una inspección exhaustiva del UAS, la estación de control (GCS) y los equipos asociados.
    \item \textbf{Análisis de Seguridad en Sitio:} Previo al inicio de cada proyecto, el Jefe de Pilotos o un piloto designado realizará un análisis en el lugar de la operación para identificar peligros específicos del entorno: obstáculos (naturales y artificiales), infraestructura crítica, áreas de posible interferencia electromagnética, condiciones de acceso y la presencia de personas ajenas a la operación.
\end{itemize}

\textbf{b) Herramienta de Evaluación de Riesgo de Vuelo (FRAT):} Antes de cada operación, el piloto remoto deberá completar el formato \textbf{F-011 Evaluación del Riesgo Operacional (Anexo 6 - Guía de Uso FRAT)}. Esta herramienta calcula el riesgo presente de la operación, asignando una puntuación de 1 a 5 a 12 factores de riesgo (ej. relacionados con el piloto, el entorno, la complejidad de la misión). Un puntaje elevado alertará al equipo sobre la necesidad de implementar medidas de mitigación adicionales antes de iniciar el vuelo.

\subsection{Fuentes de Identificación de Peligros}
\begin{longtable}{|>{\RaggedRight\arraybackslash}p{3cm}|>{\RaggedRight\arraybackslash}p{6.2cm}|>{\RaggedRight\arraybackslash}p{4.4cm}|}
\hline
\rowcolor[gray]{0.9}\textbf{Enfoque} & \textbf{Fuentes} & \textbf{Responsable} \\ \hline
\endhead
Reactivo & Reportes obligatorios y voluntarios, investigación de eventos, discrepancias de mantenimiento, quejas de clientes & Gerente SMS \\ \hline
Proactivo & FRAT, inspecciones prevuelo, análisis del sitio, auditorías, lista de verificación de la misión, observación de operaciones, cambios en procesos & Jefe de Pilotos y Gerente SMS \\ \hline
Predictivo & Tendencias de SPI, registros de vuelo y mantenimiento (por ejemplo, desgaste de motores o baterías), análisis de fatiga y estrés, cambios en normas & Gerente SMS \\ \hline
Internas & Personal, aeronaves, instalaciones, procedimientos y software & SAG \\ \hline
Externas & Aerocivil, clientes, clima, tráfico aéreo, proveedores, otros operadores & Gerente SMS \\ \hline
\end{longtable}

\textbf{Factores humanos y organizacionales.} Se miran la fatiga, el estrés, la presión del cliente, la carga de trabajo y la comunicación como fuente de peligros y se registran en el Anexo 7.

\subsection{Inicio de una Investigación y Perfil del Investigador}
\begin{itemize}
    \item \textbf{Cuándo se investiga:} siempre en eventos del Anexo 4 y en cualquier evento con riesgo rojo o amarillo alto, daño a la aeronave o repetición de un evento. En los demás casos, el Gerente SMS decide con un análisis inicial.
    \item \textbf{Quién investiga:} el Gerente SMS, que nombra por escrito al investigador designado (puede ser él mismo). Si el evento le involucra, el Gerente General nombra otra persona.
    \item \textbf{Perfil:} formación en SMS o investigación de eventos (curso del Anexo 12), conocimiento de la operación UAS, independencia frente a las personas involucradas y capacidad de entrevistar con respeto.
    \item \textbf{Contenido del informe:} hechos, cronología, causas y factores contribuyentes (incluidos los humanos y organizacionales), riesgo, acciones y responsables.
    \item \textbf{Tendencias:} cada trimestre el Gerente SMS revisa los eventos y reportes para buscar repeticiones.
\end{itemize}

\subsection{Consolidación y Registro de Peligros}

Todos los peligros identificados a través de los métodos reactivos, proactivos y predictivos serán registrados y gestionados centralizadamente en el \textbf{“Anexo 7 - Registro de Peligros de Seguridad Operacional"}. Este registro forma parte de nuestro Sistema de Recopilación y Procesamiento de Datos de Seguridad Operacional (SDCPS) y es la base para el siguiente paso del proceso de SRM.

\section{Evaluación y Mitigación de Riesgos de Seguridad Operacional}
Una vez que un peligro ha sido identificado y registrado a través de los métodos descritos en la sección anterior, es imperativo que SKYFILMS lo analice para determinar el riesgo de seguridad operacional asociado. Este proceso de gestión de riesgos (SRM) es sistemático y asegura que se tomen decisiones informadas para mantener los riesgos en un nivel tan bajo como sea razonablemente práctico (ALARP).\\

SKYFILMS aplica dos enfoques de gestión de riesgos, alineados con la naturaleza de la identificación del peligro:

\begin{enumerate}
    \item \textbf{Para Sucesos Reactivos:} Aplicado a los peligros identificados a través de notificaciones de seguridad operacional e investigaciones.
    \item \textbf{Para Operaciones Proactivas y Predictivas:} Aplicado durante la planificación de nuevas operaciones o misiones complejas, utilizando proyecciones estadísticas y auditorias internas y externas.
\end{enumerate}

\subsection{Gestión de Riesgos para Sucesos Reactivos}
Este proceso se activa cuando se recibe una notificación de seguridad (voluntaria u obligatoria) o como resultado de una investigación. La trazabilidad de estos peligros se consigna en el \textbf{Anexo 8 - Matriz de Evaluación de Riesgos}\\

\textbf{a) Índice de Riesgo de Seguridad Operacional:}\\

Para estandarizar la toma de decisiones, todo peligro identificado se evalúa utilizando una matriz de riesgo que cruza la \textbf{probabilidad} de ocurrencia de un suceso con la \textbf{gravedad} de sus consecuencias potenciales.\\

El proceso inicia con la utilización de la matriz de evaluación del riesgo, que asigna un código alfanumérico al riesgo inicial.

\begin{table}[H]
\centering
\renewcommand{\arraystretch}{1.5}
\begin{tabularx}{\textwidth}{|>{\RaggedRight\arraybackslash}p{3.4cm}|*{5}{>{\centering\arraybackslash}X|}}
\hline
\rowcolor[gray]{0.9}\textbf{Probabilidad} & \textbf{Catastrófico (A)} & \textbf{Peligroso (B)} & \textbf{Mayor (C)} & \textbf{Menor (D)} & \textbf{Insignificante (E)} \\ \hline
Frecuente (5) & \cellcolor[HTML]{FF4B4B}\textbf{5A} & \cellcolor[HTML]{FF4B4B}\textbf{5B} & \cellcolor[HTML]{FF4B4B}\textbf{5C} & \cellcolor[HTML]{FFE04D}\textbf{5D} & \cellcolor[HTML]{FFE04D}\textbf{5E} \\ \hline
Ocasional (4) & \cellcolor[HTML]{FF4B4B}\textbf{4A} & \cellcolor[HTML]{FF4B4B}\textbf{4B} & \cellcolor[HTML]{FFE04D}\textbf{4C} & \cellcolor[HTML]{FFE04D}\textbf{4D} & \cellcolor[HTML]{FFE04D}\textbf{4E} \\ \hline
Remoto (3) & \cellcolor[HTML]{FF4B4B}\textbf{3A} & \cellcolor[HTML]{FFE04D}\textbf{3B} & \cellcolor[HTML]{FFE04D}\textbf{3C} & \cellcolor[HTML]{FFE04D}\textbf{3D} & \cellcolor[HTML]{3DBE6B}\textbf{3E} \\ \hline
Improbable (2) & \cellcolor[HTML]{FFE04D}\textbf{2A} & \cellcolor[HTML]{FFE04D}\textbf{2B} & \cellcolor[HTML]{FFE04D}\textbf{2C} & \cellcolor[HTML]{3DBE6B}\textbf{2D} & \cellcolor[HTML]{3DBE6B}\textbf{2E} \\ \hline
Extremadamente improbable (1) & \cellcolor[HTML]{3DBE6B}\textbf{1A} & \cellcolor[HTML]{3DBE6B}\textbf{1B} & \cellcolor[HTML]{3DBE6B}\textbf{1C} & \cellcolor[HTML]{3DBE6B}\textbf{1D} & \cellcolor[HTML]{3DBE6B}\textbf{1E} \\ \hline
\end{tabularx}
\caption{Matriz de evaluación de riesgos: probabilidad (filas) por gravedad (columnas).}
\end{table}


Con el índice de riesgo obtenido, se utiliza la matriz de tolerabilidad para clasificar el riesgo en una de tres regiones:

\begin{itemize}
    \item \textbf{Región NO Tolerable (Roja):} El riesgo es inaceptable y las operaciones relacionadas deben cesar o no iniciarse hasta que el riesgo sea mitigado a un nivel inferior. La gestión es inmediata y requiere la intervención del Gerente SMS y la alta dirección.
    \item \textbf{Región Tolerable (Amarilla):} El riesgo es aceptable siempre que se implementen medidas de mitigación efectivas. Puede requerir una decisión de la dirección para aceptar el riesgo residual.
    \item \textbf{Región Aceptable (Verde):} El riesgo se considera aceptable y puede no requerir acciones de mitigación adicionales, aunque se mantiene bajo observación.
\end{itemize}

\begin{table}[H]
\centering
\renewcommand{\arraystretch}{1.4}
\begin{tabularx}{\textwidth}{|>{\centering\arraybackslash}p{2.6cm}|>{\centering\arraybackslash}p{3.4cm}|>{\RaggedRight\arraybackslash}X|>{\RaggedRight\arraybackslash}p{3.2cm}|}
\hline
\rowcolor[gray]{0.9}\textbf{Región} & \textbf{Índices de riesgo} & \textbf{Criterio} & \textbf{Quién decide y plazo} \\ \hline
\cellcolor[HTML]{FF4B4B}\textbf{No tolerable} & 5A, 5B, 5C, 4A, 4B, 3A & Inaceptable en las condiciones actuales; la operación se detiene. & Gerente General; respuesta inmediata \\ \hline
\cellcolor[HTML]{FFE04D}\textbf{Tolerable} & 5D, 5E, 4C, 4D, 4E, 3B, 3C, 3D, 2A, 2B, 2C & Aceptable si se aplican controles. Requiere decisión de la dirección. & Gerente SMS (ver tabla de niveles de autoridad del capítulo 2) \\ \hline
\cellcolor[HTML]{3DBE6B}\textbf{Aceptable} & 3E, 2D, 2E, 1A, 1B, 1C, 1D, 1E & Aceptable; se mantiene bajo vigilancia. & Gerente SMS; sin plazo \\ \hline
\end{tabularx}
\caption{Matriz de tolerabilidad del riesgo.}
\end{table}


La determinación de la \textbf{probabilidad y gravedad} se basa en las definiciones estandarizadas por SKYFILMS, las cuales han sido adaptadas de la guía de la OACI.

\begin{table}[H]
\centering
\renewcommand{\arraystretch}{1.4}
\begin{tabularx}{\textwidth}{|>{\RaggedRight\arraybackslash}p{4.4cm}|>{\centering\arraybackslash}X|>{\centering\arraybackslash}p{2cm}|}
\hline
\rowcolor[gray]{0.9}\textbf{Probabilidad de ocurrencia} & \textbf{Significado} & \textbf{Valor} \\ \hline
Frecuente & 70 eventos por cada 100 operaciones & \textbf{5} \\ \hline
Ocasional & 50 eventos por cada 100 operaciones & \textbf{4} \\ \hline
Remoto & 25 eventos por cada 100 operaciones & \textbf{3} \\ \hline
Improbable & 10 eventos por cada 100 operaciones & \textbf{2} \\ \hline
Extremadamente improbable & 1 evento por cada 100 operaciones & \textbf{1} \\ \hline
\end{tabularx}
\caption{Probabilidad de ocurrencia.}
\end{table}


\begin{table}[H]
\centering
\renewcommand{\arraystretch}{1.3}
\begin{tabularx}{\textwidth}{|>{\RaggedRight\arraybackslash}p{3cm}|>{\RaggedRight\arraybackslash}X|>{\centering\arraybackslash}p{1.6cm}|}
\hline
\rowcolor[gray]{0.9}\textbf{Gravedad del suceso} & \textbf{Significado} & \textbf{Valor} \\ \hline
Catastrófico & Destrucción del equipo; muertes múltiples. & \textbf{A} \\ \hline
Peligroso & Reducción importante de los márgenes de seguridad, daño físico o una carga de trabajo tal que el personal no puede desempeñar sus tareas de forma precisa y completa; lesiones graves; daños mayores al equipo. & \textbf{B} \\ \hline
Mayor & Reducción significativa de los márgenes de seguridad o de la habilidad del operador para responder a condiciones adversas por aumento de la carga de trabajo o por condiciones que impiden su eficiencia; incidente grave; lesiones a personas. & \textbf{C} \\ \hline
Menor & Interferencia; limitaciones operacionales; uso de procedimientos de emergencia; incidentes menores. & \textbf{D} \\ \hline
Insignificante & Consecuencias leves. & \textbf{E} \\ \hline
\end{tabularx}
\caption{Gravedad de las consecuencias.}
\end{table}


Este flujo de análisis, evaluación y decisión asegura un tratamiento consistente y documentado para cada peligro identificado.

\begin{table}[H]
\centering
\renewcommand{\arraystretch}{1.35}
\begin{tabularx}{\textwidth}{|c|>{\RaggedRight\arraybackslash}p{4.2cm}|>{\RaggedRight\arraybackslash}X|}
\hline
\rowcolor[gray]{0.9}\textbf{Paso} & \textbf{Actividad} & \textbf{Qué ocurre después} \\ \hline
1 & Identificación del peligro & Pasa a evaluar la probabilidad. \\ \hline
2 & Evaluación de la probabilidad (1 a 5) & Pasa a evaluar la gravedad. \\ \hline
3 & Evaluación de la gravedad (A a E) & Con el índice se ubica el riesgo en una región de la matriz de tolerabilidad. \\ \hline
4 & \cellcolor[HTML]{FF4B4B}\textbf{Intolerable} & Acciones correctivas (mitigación TRE); se evalúa de nuevo hasta que el riesgo sea tolerable. \\ \hline
5 & \cellcolor[HTML]{FFE04D}\textbf{Tolerable} & Acciones preventivas o plan de mejora (mitigación TRE); se evalúa de nuevo hasta que el riesgo sea aceptable. \\ \hline
6 & \cellcolor[HTML]{3DBE6B}\textbf{Aceptable} & Se acepta y se vigila. \\ \hline
7 & Fin & Se registra el resultado en el registro de peligros (Anexo 7). \\ \hline
\end{tabularx}
\caption{Proceso de gestión de los riesgos de seguridad operacional.}
\end{table}


\textbf{b) Mitigación y Control de los Riesgos de Seguridad Operacional:}\\

Como primera medida y antes de iniciar el proceso de mitigación de riesgos, en caso de tener mas de un riesgo se debe realizar el siguiente procedimiento de priorización:

\begin{enumerate}
    \item Agrupar los riesgos según su naturaleza (NO tolerable o Tolerable).
    \item Priorizar la mitigación de los riesgos presentes en la sección de \textbf{NO Tolerable}, y dentro de estos priorizar estos de la siguiente manera:
    \begin{enumerate}
        \item Riesgo frente a perdidas humanas. 
        \item Riesgo frente a infraestructura o tecnología.
        \item Riesgo frente al medio ambiente.
    \end{enumerate}
    \item Mitigar los riegos en ubicados en la sección \textbf{Tolerable}.
\end{enumerate}

Para todo riesgo clasificado en la región "NO Tolerable" o "Tolerable", se debe desarrollar e implementar un plan de acción para mitigarlo. Las estrategias de mitigación se centran en:

\begin{itemize}
    \item \textbf{Evitar:} Cancelar o posponer la actividad si el riesgo excede los beneficios de continuar.
    \item \textbf{Reducir:} Implementar acciones para disminuir la probabilidad de ocurrencia o la severidad de las consecuencias.
    \item \textbf{Separar la exposición:} Crear barreras o redundancias para protegerse de los efectos del peligro.
\end{itemize}

Como base principal para las defensas implementadas, SKYFILMS utiliza la metodología \textbf{TRE (Tecnología, Reglamento y Entrenamiento)}, conforme al RAC 100.715:

\begin{itemize}
    \item \textbf{Tecnología:} Uso de sistemas como por ejemplo geocercas, paracaídas balísticos, sistemas de detección y evasión (DAA), o UAS con mayor redundancia.
    \item \textbf{Reglamento:} Creación o modificación de procedimientos internos, políticas y listas de verificación documentadas en los manuales de la empresa.
    \item \textbf{Entrenamiento:} Capacitación específica al personal sobre nuevos procedimientos, simulaciones de emergencia o fortalecimiento de competencias.
\end{itemize}

Una vez aplicadas las medidas de mitigación, se realiza una nueva evaluación para determinar el \textbf{riesgo residual}. Si este se encuentra en la región ”Tolerable” aplicando las medidas propuestas en el análisis o "Aceptable", el plan de acción se considera efectivo y el peligro queda bajo control.\\

\textbf{NOTA:} El responsable de decidir si aceptar un riego ubicado en la sección \textbf{Tolerable} es el Gerente SMS.

\textbf{Tolerabilidad, autoridad y plazos.} La región de cada riesgo decide quién lo acepta y en qué plazo se mitiga; la tabla de \textit{Niveles de Autoridad para Aceptar Riesgos} está en el capítulo 2. En esa tabla, el Gerente General decide en la región roja y el Gerente SMS en la amarilla y la verde.

\textbf{Reglas para que la evaluación sea consistente.}
\begin{itemize}
    \item \textbf{Datos que faltan:} se escribe el supuesto y por qué; si falta información, se usa el peor caso creíble.
    \item \textbf{Dos evaluadores:} todo riesgo amarillo o rojo lo evalúan dos personas (el Gerente SMS y el Jefe de Pilotos o el Técnico UAS) y se justifica cualquier diferencia.
    \item \textbf{Revisión periódica:} se vuelven a evaluar los riesgos cada seis meses y siempre que haya un cambio o un evento.
    \item \textbf{Riesgo agregado:} el SAG revisa cada trimestre si varios riesgos pequeños, juntos, hacen un riesgo mayor (por ejemplo, vuelo nocturno más zona urbana más FPV).
    \item \textbf{Visibilidad:} el Gerente SMS entrega al Gerente General la lista de riesgos amarillos y rojos abiertos cada vez que se reúne el SRB.
    \item \textbf{Controles seguros:} antes de aprobar un control, se pregunta si crea un peligro nuevo (por ejemplo, un cambio de procedimiento que aumente la carga de trabajo).
\end{itemize}

\textbf{Verificación de eficacia de cada control.} Se define para cada control: método (revisar registros, repetir la inspección, observar un vuelo o mirar el SPI), frecuencia (mensual, semestral o anual, según el riesgo), responsable, registro (Anexo 8 y Anexo 11) y criterio de éxito. El resultado vuelve al registro de peligros (Anexo 7). Los controles a cargo de un proveedor se verifican con su informe y con una visita o auditoría anual.



% Capítulo 4: Aseguramiento de la Seguridad Operacional
\newgeometry{top=4.2cm, bottom=2.5cm, left=2cm, right=2cm, headheight=60pt, headsep=1.5cm}
\chapter{Aseguramiento de la Seguridad Operacional}
\label{chap:Aseguramiento de la Seguridad Operacional}
El aseguramiento de la seguridad operacional constituye el mecanismo mediante el cual SKYFILMS verifica de manera sistemática la eficacia de su Sistema de Gestión de la Seguridad Operacional (SMS). Esta sección establece los procesos que permiten evaluar el desempeño de la organización frente a sus objetivos de seguridad, garantizar la gestión controlada de los cambios y promover la mejora continua.\\

El proceso general de gestión del rendimiento en seguridad operacional y su relación con los sistemas de recopilación y procesamiento de datos sobre seguridad operacional (SDCPS) se resume en la siguiente tabla:

\begin{table}[H]
\centering
\renewcommand{\arraystretch}{1.35}
\begin{tabularx}{\textwidth}{|>{\RaggedRight\arraybackslash}p{4.6cm}|>{\RaggedRight\arraybackslash}X|}
\hline
\rowcolor[gray]{0.9}\textbf{Elemento del SDCPS} & \textbf{Qué hace} \\ \hline
\rowcolor[HTML]{FFE04D}\textbf{Rendimiento del SMS} & Es el centro del sistema: todos los demás elementos lo alimentan y se retroalimentan con él. \\ \hline
Definición y actualización de SPI & Define los indicadores de seguridad operacional (Anexo 9, tabla A) y los actualiza. \\ \hline
Definición y actualización de KPI & Define los indicadores internos de gestión (Anexo 9, tabla B). \\ \hline
Medición del rendimiento & Mide cada mes los ciclos de vuelo y los eventos y calcula las tasas. \\ \hline
Análisis del SMS & El SAG analiza las tendencias, causas y alertas. \\ \hline
Promoción del SMS & Comunica los resultados y las lecciones aprendidas al personal. \\ \hline
Identificación de medidas necesarias & Define las acciones y planes de mejora que se derivan de los resultados. \\ \hline
\end{tabularx}
\caption{Sistemas de recopilación y procesamiento de datos sobre seguridad operacional (SDCPS).}
\end{table}


\section{Observación y Medición del Rendimiento}
SKYFILMS implementará un sistema estructurado de observación y medición del rendimiento en seguridad operacional, en cumplimiento del RAC 219, la Circular MAUT-1.0-22-005 de Aerocivil y las directrices del Doc 9859 de la OACI.

\begin{itemize}
    \item La organización establecerá y mantendrá Indicadores de Desempeño en Seguridad Operacional (SPI), los cuales estarán directamente vinculados con los objetivos de seguridad de SKYFILMS y con los niveles aceptables de seguridad operacional aprobados por la Autoridad Aeronáutica.
    \item La recopilación, análisis y reporte de los SPI se realizará de forma periódica, garantizando la trazabilidad de la información y permitiendo identificar tendencias, desviaciones y oportunidades de mejora.
    \item Los resultados serán evaluados anualmente, y servirán de base para la toma de decisiones en materia de mitigación de riesgos, asignación de recursos y definición de acciones correctivas y preventivas.
    \item Toda la información asociada al monitoreo de SPI se documentará en formatos oficiales de la organización y será consolidada en un anexo específico al presente manual \textbf{(Anexo 9 - Indicadores de Rendimiento de Seguridad Operacional SPI's)}, que incluirá la descripción detallada de cada indicador, su fórmula, metas, plan de acción y mecanismos de evaluación.
\end{itemize}



Los SPI de Skyfilms son los once indicadores que la Aerocivil concertó para explotadores UAS (circular MAUT-1.0-22-005). Sus fórmulas (eventos por cada 1\,000 ciclos de vuelo), metas, niveles de alerta y forma de entrega están en el Anexo 9, y su relación con los objetivos de seguridad operacional en el Anexo 1. El Gerente SMS anota los datos cada mes, los presenta al SAG cada trimestre y el SRB los revisa al menos una vez al año. Cada año, antes del 30 de marzo, se envía el archivo SPI a la Aerocivil.

\section{Gestión del Cambio}
SKYFILMS reconoce que la introducción de cambios en el funcionamiento de la organización puede generar nuevos riesgos o modificar los existentes, debido a esto a continuación se listan los principales desencadenantes que activan el proceso de gestión del cambio:

\begin{enumerate}
    \item Cambios las operaciones.
    \item Cambios o inclusión de equipos.
    \item Cambios o inclusión de personal.
    \item Cambios o inclusión de procedimientos.
\end{enumerate}

 Por lo tanto, se aplicará un proceso formal de gestión del cambio \textbf{(Anexo 10 - Procedimiento de Gestión del Cambio)} para garantizar que toda modificación sea evaluada y aprobada antes de su implementación.


\textbf{Desencadenantes, responsables, plazos y seguimiento:} están en el Anexo 10. Incluyen la participación de las personas afectadas, la aprobación antes de aplicar el cambio, la revisión de las evaluaciones de riesgo ya hechas, el análisis de riesgos transitorios y acumulados, la comunicación y el seguimiento a los 30 y 90 días.

\section{Programa de Auditoria Interna}

\subsection{Propósito}

SKYFILMS establecerá y mantendrá un Programa de Auditoría Interna como herramienta sistemática, independiente y documentada para evaluar:

\begin{itemize}
    \item La conformidad del Sistema de Gestión de la Seguridad Operacional (SMS) con la normativa aplicable.

    \item El cumplimiento de las políticas, procedimientos y controles establecidos en el presente Manual.

    \item La eficacia de los controles de mitigación implementados.

    \item El desempeño de los Indicadores de Seguridad Operacional (SPI).

    \item La madurez y efectividad global del SMS.
\end{itemize}

El programa constituye un elemento fundamental del proceso de aseguramiento de la seguridad operacional y soporte a la mejora continua.

\subsection{Alcance}
El Programa de Auditoría Interna aplicará a:

\begin{itemize}
    \item Operaciones UAS.

    \item Gestión técnica y mantenimiento.

    \item Gestión documental del SMS.

    \item Gestión de riesgos.

    \item Gestión del cambio.

    \item Gestión de reportes voluntarios.

    \item Cumplimiento normativo (RAC 100, RAC 219 y demás regulación aplicable).

    \item Gestión de instrucción y competencia del personal.

    \item Cumplimiento de procedimientos operacionales.
\end{itemize}

\subsection{Periodicidad}
El programa anual de auditorías será aprobado por el Gerente SMS y el Gerente General, de la cual realizará como mínimo una auditoría interna \textbf{anual} completa del Sistema de Gestión de Seguridad Operacional SMS, y ademas podrán programarse auditorías parciales o específicas cuando:

\begin{itemize}
    \item Se identifiquen desviaciones significativas en los SPI.

    \item Existan hallazgos recurrentes.

    \item Se implementen cambios operacionales relevantes.

    \item La Dirección lo considere necesario.
\end{itemize}

\subsection{Metodologia}
Las auditorías internas se desarrollarán bajo las siguientes etapas:

\begin{enumerate}
    \item Planificación de la auditoría.

    \item Revisión documental.

    \item Entrevistas y verificación en campo.

    \item Evaluación de cumplimiento.

    \item Identificación de hallazgos.

    \item Emisión de informe formal.

    \item Seguimiento a acciones correctivas.

    \item Verificación de eficacia.
\end{enumerate}

\subsection{Calendario, Independencia y Competencia del Auditor}
\begin{itemize}
    \item \textbf{Calendario:} el programa anual se fija en enero. Se audita una vez al año cada área (operaciones, mantenimiento, documentos, riesgos, capacitación) y el SMS completo. Las fechas se registran en el programa y se avisan con 15 días de anticipación.
    \item \textbf{Independencia:} nadie audita su propio trabajo. El Gerente SMS audita operaciones y mantenimiento; al SMS lo audita el Jefe de Pilotos con apoyo de un auditor externo o de una persona sin relación con el área auditada.
    \item \textbf{Competencia del auditor:} curso de SMS y conocimiento del RAC 100 y el RAC 219; el curso de personal clave del Anexo 12 cumple este requisito.
    \item \textbf{Contratistas y proveedores:} se audita al proveedor de mantenimiento (cuando exista) y la ejecución de sus controles, al menos una vez al año.
    \item \textbf{Acceso al Ejecutivo Responsable:} el responsable de la auditoría entrega el informe directamente al Gerente General y puede pedirle una reunión en cualquier momento.
    \item \textbf{Informes y cierre:} el informe sale dentro de 15 días hábiles después de la auditoría; los hallazgos se cierran según el Plan de Mejora Continua (Anexo 11) y los registros se guardan cinco años.
\end{itemize}

\subsection{Clasificación de Hallazgos}

Los hallazgos derivados de auditoría se clasificarán como:\\

\textbf{Hallazgo Mayor:}

\begin{itemize}
    \item Incumplimiento que afecta directamente la seguridad operacional.

    \item Ausencia o falla crítica de un control.

    \item Incumplimiento normativo significativo.
\end{itemize}

\textbf{Hallazgo Menor:}

\begin{itemize}
    \item Desviación puntual sin impacto inmediato en la seguridad.

    \item Oportunidad de mejora.
\end{itemize}

\textbf{Observación:}

\begin{itemize}
    \item Recomendación preventiva sin incumplimiento directo.
\end{itemize}

\subsection{Gestión de Hallazgos}
Todos los hallazgos deberán:

\begin{itemize}
    \item Registrarse en el Formato de Plan de Mejora Continua.

    \item Definir acción correctiva o preventiva.

    \item Asignar responsable.

    \item Establecer plazo de cierre.

    \item Verificar eficacia antes de cierre formal.
\end{itemize}

El Gerente SMS realizará seguimiento periódico hasta su cierre.

\subsection{Informe de Auditoria}
Cada auditoría generará un informe formal que contendrá:

\begin{itemize}
    \item Alcance

    \item Objetivo

    \item Criterios utilizados

    \item Metodología aplicada

    \item Hallazgos identificados

    \item Nivel de cumplimiento

    \item Conclusión general sobre el estado del SMS
\end{itemize}

El informe será presentado a la Gerencia para su conocimiento y aprobación de acciones.

\subsection{Conservación de Registros}
Todos los registros asociados al Programa de Auditoría Interna serán conservados conforme al sistema de gestión documental del SMS por un \textbf{período mínimo de cinco (5) años}.

\subsection{Integración con el Sistema SMS}

El Programa de Auditoría Interna alimenta directamente:

\begin{enumerate}
    \item La medición del desempeño del SMS.

    \item La revisión de los SPI.

    \item La gestión del cambio.

    \item El Plan de Mejora Continua.

\end{enumerate}

Lo que constituye un mecanismo estructural para asegurar la efectividad y el seguimiento de los datos utilizados para la toma de decisiones de gestión y asegurar la mejora continua sostenida del Sistema de Gestión de Seguridad Operacional de SKYFILMS.

\section{Mejora Continua}
La seguridad operacional es un proceso dinámico que requiere de revisión y optimización permanentes. SKYFILMS se compromete a mantener un ciclo de mejora continua del SMS, garantizando que los resultados obtenidos conduzcan a un fortalecimiento progresivo del sistema.

\begin{itemize}
    \item Se realizarán auditorías internas de seguridad operacional con una periodicidad mínima \textbf{anual}, así como revisiones por la dirección al menos una vez por año.
    \item Los hallazgos derivados de auditorías, inspecciones, encuestas de seguridad, reportes voluntarios y el análisis de los SPI serán integrados en un Plan de Mejora Continua \textbf{(Anexo 11 - Plan de Mejora Continua)}, que definirá acciones correctivas, responsables, plazos y mecanismos de verificación.
    \item Toda acción correctiva o preventiva deberá estar sustentada en datos objetivos y será verificada para confirmar su eficacia antes de ser cerrada.
    \item La evolución del desempeño de seguridad operacional, las lecciones aprendidas y las mejoras implementadas serán comunicadas de manera transparente a todo el personal, fortaleciendo la cultura de seguridad en la organización.
    \item Como responsable principal esta el Gerente SMS quien sera el encargado de liderar las auditorias y revisiones.
\end{itemize}





% Capítulo 5: Promoción de la Seguridad Operacional
\newgeometry{top=4.2cm, bottom=2.5cm, left=2cm, right=2cm, headheight=60pt, headsep=1.5cm}
\chapter{Promoción de la Seguridad Operacional}
\label{chap:Promoción de la Seguridad Operacional}
La promoción de la seguridad operacional en SKYFILMS constituye el conjunto de actividades orientadas a fortalecer la \textbf{cultura de seguridad}, garantizar la \textbf{competencia del personal} y asegurar una \textbf{comunicación efectiva y transversal} de la información en todos los niveles de la organización.\\

El éxito del SMS depende de que cada miembro de la organización comprenda sus responsabilidades, reciba formación adecuada y disponga de canales de comunicación que fomenten la transparencia y el aprendizaje colectivo.

\section{Instrucción y Educación}
SKYFILMS establece un \textbf{programa integral de instrucción y educación en seguridad operacional}, con el objetivo de asegurar que todo el personal involucrado en la operación de UAS cuente con las competencias necesarias para desempeñar sus funciones de manera segura, eficiente y en conformidad con la normativa vigente.\\

\textbf{Principios rectores:}

\begin{itemize}
    \item La instrucción en seguridad operacional será obligatoria para todo el personal que participa directa o indirectamente en las operaciones de UAS.
    \item La formación incluirá \textbf{inducción inicial, entrenamiento recurrente y capacitaciones específicas} asociadas a riesgos identificados en la operación.
    \item El contenido de los programas de capacitación será coherente con los objetivos de seguridad operacional de SKYFILMS, los SPI definidos, y la normativa aplicable (RAC 100 y RAC 219).
    \item Las temáticas de instrucción incluirán de manera obligatoria las \textbf{políticas y procedimientos del explotador UAS}, así como contenidos de \textbf{Factores Humanos (FFHH)} y otras buenas prácticas orientadas a la prevención de errores y a la promoción de la seguridad de la operación.
    \item Se fomentará la participación en \textbf{simulacros de emergencias operacionales y ejercicios prácticos}, con el fin de fortalecer la respuesta del personal ante situaciones críticas.
    \item Los registros de asistencia y evidencias de capacitación se conservarán conforme al sistema de gestión documental del SMS.
\end{itemize}

\textbf{Objetivos específicos de la instrucción:}

\begin{enumerate}
    \item Garantizar el conocimiento y aplicación de los procedimientos de seguridad operacional establecidos.
    \item Promover la \textbf{cultura justa}, asegurando que los empleados comprendan la importancia del reporte voluntario de peligros sin temor a represalias.
    \item Reforzar las competencias técnicas de pilotos, observadores y personal de apoyo en escenarios de riesgo operacional identificados.
    \item Asegurar que la instrucción evolucione de acuerdo con las lecciones aprendidas, los hallazgos de auditorías y las tendencias de seguridad observadas en los SPI.
    \item Asegurar la comprensión de los \textbf{Factores Humanos (FFHH)} y la aplicación de \textbf{buenas prácticas operacionales} como elementos clave en la mejora de la seguridad operacional.
\end{enumerate}

\subsection{Competencias del Personal Clave en Seguridad Operacional}

Con el fin de garantizar la eficacia del Sistema de Gestión de la Seguridad Operacional (SMS), SKYFILMS establecerá y mantendrá definidas las competencias mínimas requeridas para el personal clave en materia de seguridad operacional, asegurando que el programa de instrucción cubra de manera integral las funciones individuales, responsabilidades y obligaciones de rendición de cuentas dentro del sistema.\\

\subsubsection{Competencias del Personal Clave de Seguridad Operacional}

    El personal con responsabilidades directas dentro del SMS deberá demostrar, como mínimo, las siguientes competencias:
    
    \begin{enumerate}
        \item \textbf{Responsable del SMS:}
            \begin{itemize}
                \item Conocimiento sólido de la normativa aplicable (RAC 100, RAC 219 y disposiciones complementarias).

                \item Dominio de los cuatro componentes del SMS.

                \item Capacidad para liderar procesos de identificación de peligros y gestión del riesgo.

                \item Competencia en análisis de datos de seguridad e interpretación de SPI.

                \item Habilidades en auditoría interna y mejora continua.

                \item Capacidad de comunicación efectiva y liderazgo en cultura justa.
            \end{itemize}
        \item \textbf{Personal Operativo (Pilotos UAS, Observadores, Personal de Apoyo):}
            \begin{itemize}
                \item Conocimiento de la Política y Objetivos de Seguridad Operacional.

                \item Comprensión de sus responsabilidades individuales dentro del SMS.

                \item Competencia en procedimientos operacionales normales y de emergencia.

                \item Capacidad para identificar peligros y reportarlos oportunamente.

                \item Aplicación de principios de gestión del riesgo en la toma de decisiones operacionales.
            \end{itemize}
    \end{enumerate}
    
\subsubsection{Cobertura del Programa de Instrucción SMS}

    El programa de instrucción en seguridad operacional garantizará que todo el personal:
        \begin{itemize}
            \item Comprenda cómo funciona el SMS dentro de SKYFILMS.

            \item Conozca los canales formales de reporte de peligros y sucesos.

            \item Identifique claramente sus roles y responsabilidades individuales.

            \item Entienda su obligación de rendición de cuentas frente al cumplimiento de procedimientos.

            \item Reconozca la importancia de la cultura justa y del reporte voluntario.

            \item Comprenda la relación entre objetivos de seguridad, SPI y desempeño operacional.
        \end{itemize}\\

    La instrucción incluirá, como mínimo:

        \begin{itemize}
            \item Estructura y componentes del SMS.

            \item Gestión del riesgo operacional (identificación de peligros, evaluación y mitigación).

            \item Uso del sistema de reporte.

            \item Análisis básico de eventos y lecciones aprendidas.

            \item Responsabilidad individual en la prevención de incidentes.

            \item Factores humanos aplicables a operaciones UAS.

            \item Buenas prácticas operacionales y disciplina operacional.
        \end{itemize}

\textbf{NOTA:} Para la implementación se seguirá el \textbf{Anexo 12 - Programa de Instrucción SMS}.
\subsubsection{Evaluación de Competencias}

    La organización verificará el cumplimiento de las competencias mediante:
        \begin{itemize}
            \item Evaluaciones teóricas en instrucción inicial y recurrente.

            \item Participación en simulacros y ejercicios prácticos.

            \item Revisión de desempeño operacional.

            \item Seguimiento de participación en actividades del SMS.

            \item Retroalimentación derivada de auditorías internas y análisis de SPI.
        \end{itemize}



\section{Comunicación de la Seguridad Operacional}
SKYFILMS fomentará una \textbf{comunicación clara, efectiva y bidireccional} en materia de seguridad operacional, entendida como un elemento esencial para consolidar una cultura organizacional positiva y orientada a la prevención de riesgos.\\

\textbf{Lineamientos de comunicación:}

\begin{itemize}
    \item Establecimiento de un \textbf{sistema de reporte voluntario y confidencial}, accesible para todos los colaboradores, como herramienta clave para la identificación temprana de peligros.
    \item Difusión periódica de \textbf{boletines de seguridad operacional}, que incluyan estadísticas de desempeño, análisis de incidentes, lecciones aprendidas y mejores prácticas.
    \item Realización de \textbf{reuniones de seguridad operacional (Safety Meetings)} lideradas por el Gerente de Seguridad Operacional, con el fin de presentar avances, retroalimentar al personal y reforzar el compromiso de la dirección.
    \item Implementación de \textbf{canales digitales} (intranet, correos institucionales, grupos de trabajo) para garantizar el acceso oportuno a la información de seguridad.
    \item Comunicación inmediata de \textbf{alertas de seguridad operacional} cuando se identifiquen tendencias de riesgo que requieran acciones preventivas urgentes.
    \item Coordinación con terceros, autoridades aeronáuticas y clientes para compartir información relevante que impacte la seguridad de las operaciones.

\end{itemize}

\textbf{NOTA:} En caso de tener información critica a comunicar a las personas involucradas en la operación UAS, el Gerente de Segurar Operacional usara, de los métodos y canales anteriormente mencionados, el método mas eficaz según el contexto.\\

\textbf{Objetivos específicos de la comunicación:}

\begin{enumerate}
    \item Mantener al personal informado sobre los \textbf{resultados del desempeño del SMS} y las acciones derivadas de auditorías, inspecciones y análisis de SPI.
    \item Facilitar la retroalimentación desde todos los niveles de la organización, garantizando que las inquietudes en materia de seguridad sean consideradas en la toma de decisiones.
    \item Generar confianza y compromiso en los colaboradores mediante la transparencia en el manejo de la información.
    \item Promover el aprendizaje organizacional a través de la divulgación de lecciones aprendidas propias y de otros operadores UAS.
\end{enumerate}




% Apendice: Anexos
\chapter*{Apéndices}          % Sección sin número
\addcontentsline{toc}{chapter}{Apéndices} % Lo añade al índice

\section*{A. Anexos}
\addcontentsline{toc}{section}{A. Anexos}

\begin{itemize}
    \item \textbf{Anexo 1:} Objetivos de Seguridad Operacional KPI's.
    
    \item \textbf{Anexo 2:} Contactos de Emergencia.
    
    \item \textbf{Anexo 3:} Procedimiento de Control de Documentos y Registros.

    \item \textbf{Anexo 4:} Listado Reportes Obligatorios M.O.R. 

    \item \textbf{Anexo 5:} Guía de Uso Formato de Notificación Peligros de Seguridad Operacional.

    \item \textbf{Anexo 6:} Guía de Uso Evaluación de Riesgos de Vuelo FRAT.

    \item \textbf{Anexo 7:} Registro de Peligros de Seguridad Operacional.

    \item \textbf{Anexo 8:} Matriz de Evaluación de Riesgos (Metodo Reactivo).

    \item \textbf{Anexo 9:} Indicadores de Rendimiento de Seguridad Operacional SPI's.

    \item \textbf{Anexo 10:} Procedimiento de Gestión del Cambio.

    \item \textbf{Anexo 11:} Plan de Mejora Continua.

    \item \textbf{Anexo 12:} Programa de Instrucción SMS.
\end{itemize}

\section*{B. Documentación de Referencia}
\addcontentsline{toc}{section}{B. Documentación de Referencia}

\begin{itemize}
    \item Doc 9859 - Manual de Gestión de Seguridad Operacional
    \item RAC 219 - Implementación del Sistema SMS
    \item RAC 100 - Operación de Sistemas De Aeronaves No Tripuladas UAS 
    \item 100-MAUT-5.0-22-017 CI-Para la Implementación del SMS en explotadores UAS
    \item 219-MAUT-1.0-22-005 CI Definicion de Indicadores de rendimiento en Materia de Seguridad Operacional
    \item Circular GIVC-1.0-22-015 Respuesta ante emergencias
\end{itemize}

\end{document}


